%%%%%%%%%%%%%%%%%%%%%%%%%%%%%%%%%%%%%%%%%%%%%%%%%%%%%%%%%%%%%%%%%%%%%%%%%%%%%%%
% Harvard Academic Notes - 통합 마스터 템플릿
% 모든 강의 노트에 적용되는 통일된 스타일
% 버전: 2.1 - 가독성 개선 (선택적 최적화)
% 최종 수정일: 2025-11-17
%%%%%%%%%%%%%%%%%%%%%%%%%%%%%%%%%%%%%%%%%%%%%%%%%%%%%%%%%%%%%%%%%%%%%%%%%%%%%%%

\documentclass[11pt,a4paper]{article}

%========================================================================================
% 기본 패키지
%========================================================================================

% --- 한국어 지원 ---
\usepackage{kotex}

% --- 페이지 레이아웃 ---
\usepackage[top=20mm, bottom=20mm, left=20mm, right=18mm]{geometry}
\usepackage{setspace}
\onehalfspacing                      % 1.5배 줄간격
\setlength{\parskip}{0.5em}          % 문단 간격
\setlength{\parindent}{0pt}          % 들여쓰기 없음

% --- 표 관련 ---
\usepackage{booktabs}              % 고품질 표
\usepackage{tabularx}              % 자동 너비 조절 표
\usepackage{array}                 % 표 컬럼 확장
\usepackage{longtable}             % 여러 페이지 표
\renewcommand{\arraystretch}{1.1}  % 표 행간 조절

%========================================================================================
% 헤더 및 푸터
%========================================================================================

\usepackage{fancyhdr}
\pagestyle{fancy}
\fancyhf{}
\fancyhead[L]{\small\textit{FIN-571: Financial Institutions & Monetary Policy}}
\fancyhead[R]{\small\textit{Module 01}}
\fancyfoot[C]{\thepage}
\renewcommand{\headrulewidth}{0.5pt}
\renewcommand{\footrulewidth}{0.3pt}

% 첫 페이지는 헤더 없음
\fancypagestyle{firstpage}{
    \fancyhf{}
    \fancyfoot[C]{\thepage}
    \renewcommand{\headrulewidth}{0pt}
}

%========================================================================================
% 색상 정의 (파스텔 톤 + 다크모드 호환)
%========================================================================================

\usepackage[dvipsnames]{xcolor}

% 밝은 배경용 파스텔 색상
\definecolor{lightblue}{RGB}{220, 235, 255}      % 부드러운 파랑
\definecolor{lightgreen}{RGB}{220, 255, 235}     % 부드러운 초록
\definecolor{lightyellow}{RGB}{255, 250, 220}    % 부드러운 노랑
\definecolor{lightpurple}{RGB}{240, 230, 255}    % 부드러운 보라
\definecolor{lightgray}{gray}{0.95}              % 밝은 회색
\definecolor{lightpink}{RGB}{255, 235, 245}      % 부드러운 핑크
\definecolor{boxgray}{gray}{0.95}
\definecolor{boxblue}{rgb}{0.9, 0.95, 1.0}
\definecolor{boxred}{rgb}{1.0, 0.95, 0.95}

% 진한 색상 (테두리/제목용)
\definecolor{darkblue}{RGB}{50, 80, 150}
\definecolor{darkgreen}{RGB}{40, 120, 70}
\definecolor{darkorange}{RGB}{200, 100, 30}
\definecolor{darkpurple}{RGB}{100, 60, 150}

%========================================================================================
% 박스 환경 (tcolorbox) - 6가지 타입
%========================================================================================

\usepackage[most]{tcolorbox}
\tcbuselibrary{skins, breakable}

% 1. 개요 박스 (강의 시작 부분)
\newtcolorbox{overviewbox}[1][]{
    enhanced,
    colback=lightpurple,
    colframe=darkpurple,
    fonttitle=\bfseries\large,
    title=📚 강의 개요,
    arc=3mm,
    boxrule=1pt,
    left=8pt,
    right=8pt,
    top=8pt,
    bottom=8pt,
    breakable,
    #1
}

% 2. 요약 박스
\newtcolorbox{summarybox}[1][]{
    enhanced,
    colback=lightblue,
    colframe=darkblue,
    fonttitle=\bfseries,
    title=📝 핵심 요약,
    arc=2mm,
    boxrule=0.7pt,
    left=6pt,
    right=6pt,
    top=6pt,
    bottom=6pt,
    breakable,
    #1
}

% 3. 핵심 정보 박스
\newtcolorbox{infobox}[1][]{
    enhanced,
    colback=lightgreen,
    colframe=darkgreen,
    fonttitle=\bfseries,
    title=💡 핵심 정보,
    arc=2mm,
    boxrule=0.7pt,
    left=6pt,
    right=6pt,
    top=6pt,
    bottom=6pt,
    breakable,
    #1
}

% 4. 주의사항 박스
\newtcolorbox{warningbox}[1][]{
    enhanced,
    colback=lightyellow,
    colframe=darkorange,
    fonttitle=\bfseries,
    title=⚠️ 주의사항,
    arc=2mm,
    boxrule=0.7pt,
    left=6pt,
    right=6pt,
    top=6pt,
    bottom=6pt,
    breakable,
    #1
}

% 5. 예제 박스
\newtcolorbox{examplebox}[1][]{
    enhanced,
    colback=lightgray,
    colframe=black!60,
    fonttitle=\bfseries,
    title=📖 예제,
    arc=2mm,
    boxrule=0.7pt,
    left=6pt,
    right=6pt,
    top=6pt,
    bottom=6pt,
    breakable,
    #1
}

% 6. 정의 박스
\newtcolorbox{definitionbox}[1][]{
    enhanced,
    colback=lightpink,
    colframe=purple!70!black,
    fonttitle=\bfseries,
    title=📌 정의,
    arc=2mm,
    boxrule=0.7pt,
    left=6pt,
    right=6pt,
    top=6pt,
    bottom=6pt,
    breakable,
    #1
}

% 7. 중요 박스 (importantbox - warningbox와 유사)
\newtcolorbox{importantbox}[1][]{
    enhanced,
    colback=boxred,
    colframe=red!70!black,
    fonttitle=\bfseries,
    title=⚠️ 매우 중요,
    arc=2mm,
    boxrule=0.7pt,
    left=6pt,
    right=6pt,
    top=6pt,
    bottom=6pt,
    breakable,
    #1
}

% 8. cautionbox (warningbox와 동일)
\let\cautionbox\warningbox
\let\endcautionbox\endwarningbox

\tcbset{
    summarybox/.style={enhanced, colback=lightblue, colframe=darkblue, fonttitle=\bfseries, title=핵심 요약, arc=2mm, boxrule=0.7pt, breakable},
    warningbox/.style={enhanced, colback=lightyellow, colframe=darkorange, fonttitle=\bfseries, title=주의, arc=2mm, boxrule=0.7pt, breakable},
    examplebox/.style={enhanced, colback=lightgray, colframe=black!60, fonttitle=\bfseries, title=예제, arc=2mm, boxrule=0.7pt, breakable},
    definitionbox/.style={enhanced, colback=lightpink, colframe=purple!70!black, fonttitle=\bfseries, title=정의, arc=2mm, boxrule=0.7pt, breakable},
    importantbox/.style={enhanced, colback=boxred, colframe=red!70!black, fonttitle=\bfseries, title=매우 중요, arc=2mm, boxrule=0.7pt, breakable},
    infobox/.style={enhanced, colback=lightgreen, colframe=darkgreen, fonttitle=\bfseries, title=핵심 정보, arc=2mm, boxrule=0.7pt, breakable},
    conceptbox/.style={enhanced, colback=lightgreen, colframe=darkgreen, fonttitle=\bfseries, title=개념, arc=2mm, boxrule=0.7pt, breakable},
    analogybox/.style={enhanced, colback=green!5!white, colframe=green!60!black, fonttitle=\bfseries, title=비유, arc=2mm, boxrule=0.7pt, breakable},
    overviewbox/.style={enhanced, colback=lightpurple, colframe=darkpurple, fonttitle=\bfseries\large, title=강의 개요, arc=3mm, boxrule=1pt, breakable},
    cautionbox/.style={enhanced, colback=lightyellow, colframe=darkorange, fonttitle=\bfseries, title=주의, arc=2mm, boxrule=0.7pt, breakable},
}

%========================================================================================
% 코드 블록 설정 (밝은 배경)
%========================================================================================

\usepackage{listings}

\definecolor{codegray}{rgb}{0.5,0.5,0.5}
\definecolor{codepurple}{rgb}{0.58,0,0.82}
\definecolor{backcolour}{rgb}{0.95,0.95,0.95}

\lstset{
    basicstyle=\ttfamily\small,
    backgroundcolor=\color{lightgray},
    keywordstyle=\color{darkblue}\bfseries,
    commentstyle=\color{darkgreen}\itshape,
    stringstyle=\color{purple!80!black},
    numberstyle=\tiny\color{black!60},
    numbers=left,
    numbersep=8pt,
    breaklines=true,
    breakatwhitespace=false,
    frame=single,
    frameround=tttt,
    rulecolor=\color{black!30},
    captionpos=b,
    showstringspaces=false,
    tabsize=2,
    xleftmargin=15pt,
    xrightmargin=5pt,
    escapeinside={\%*}{*)}
}

% Python 코드 스타일
\lstdefinestyle{pythonstyle}{
    language=Python,
    morekeywords={self, True, False, None},
}

% SQL 코드 스타일
\lstdefinestyle{sqlstyle}{
    language=SQL,
    morekeywords={SELECT, FROM, WHERE, JOIN, GROUP, BY, ORDER, HAVING},
}

%========================================================================================
% 목차 스타일링
%========================================================================================

\usepackage{tocloft}
\renewcommand{\cftsecleader}{\cftdotfill{\cftdotsep}}
\setlength{\cftbeforesecskip}{0.4em}
\renewcommand{\cftsecfont}{\bfseries}
\renewcommand{\cftsubsecfont}{\normalfont}

%========================================================================================
% 표 및 그림
%========================================================================================

\usepackage{graphicx}              % 이미지
\usepackage{adjustbox}             % 표/박스 크기 조절

% 표 캡션 스타일
\usepackage{caption}
\captionsetup[table]{
    labelfont=bf,
    textfont=it,
    skip=5pt
}
\captionsetup[figure]{
    labelfont=bf,
    textfont=it,
    skip=5pt
}

%========================================================================================
% 수학
%========================================================================================

\usepackage{amsmath, amssymb, amsthm}

% 정리 환경
\theoremstyle{definition}
\newtheorem{theorem}{정리}[section]
\newtheorem{lemma}[theorem]{보조정리}
\newtheorem{proposition}[theorem]{명제}
\newtheorem{corollary}[theorem]{따름정리}
\newtheorem{definition}{정의}[section]
\newtheorem{example}{예제}[section]

%========================================================================================
% 하이퍼링크
%========================================================================================

\usepackage[
    colorlinks=true,
    linkcolor=blue!80!black,
    urlcolor=blue!80!black,
    citecolor=green!60!black,
    bookmarks=true,
    bookmarksnumbered=true,
    pdfborder={0 0 0}
]{hyperref}

% PDF 메타데이터는 각 문서에서 설정
\hypersetup{
    pdftitle={FIN-571: Financial Institutions & Monetary Policy - Module 01},
    pdfauthor={강의 노트},
    pdfsubject={Academic Notes}
}

%========================================================================================
% 기타 유용한 패키지
%========================================================================================

\usepackage{enumitem}              % 리스트 커스터마이징
\setlist{nosep, leftmargin=*, itemsep=0.3em}

\usepackage{microtype}             % 타이포그래피 개선
\usepackage{footnote}              % 각주 개선
\usepackage{url}                   % URL 줄바꿈
\urlstyle{same}

%========================================================================================
% 사용자 정의 명령어
%========================================================================================

% 강조 텍스트
\newcommand{\important}[1]{\textbf{\textcolor{red!70!black}{#1}}}
\newcommand{\keyword}[1]{\textbf{#1}}
\newcommand{\term}[1]{\textit{#1}}
\newcommand{\code}[1]{\texttt{#1}}

% 용어 설명 (인라인)
\newcommand{\defterm}[2]{\textbf{#1}\footnote{#2}}

% 섹션 시작 전 페이지 분리
\newcommand{\newsection}[1]{\newpage\section{#1}}

%========================================================================================
% 문서 제목 스타일
%========================================================================================

\usepackage{titling}
\pretitle{\begin{center}\LARGE\bfseries}
\posttitle{\par\end{center}\vskip 0.5em}
\preauthor{\begin{center}\large}
\postauthor{\end{center}}
\predate{\begin{center}\large}
\postdate{\par\end{center}}

%========================================================================================
% 섹션 제목 간격
%========================================================================================

\usepackage{titlesec}
\titlespacing*{\section}{0pt}{1.5em}{0.8em}
\titlespacing*{\subsection}{0pt}{1.2em}{0.6em}
\titlespacing*{\subsubsection}{0pt}{1em}{0.5em}

%========================================================================================
% 메타 정보 박스 명령어
%========================================================================================

\newcommand{\metainfo}[4]{
\begin{tcolorbox}[
    colback=lightpurple,
    colframe=darkpurple,
    boxrule=1pt,
    arc=2mm,
    left=10pt,
    right=10pt,
    top=8pt,
    bottom=8pt
]
\begin{tabular}{@{}rl@{}}
▣ \textbf{강의명:} & #1 \\[0.3em]
▣ \textbf{주차:} & #2 \\[0.3em]
▣ \textbf{교수명:} & #3 \\[0.3em]
▣ \textbf{목적:} & \begin{minipage}[t]{0.75\textwidth}#4\end{minipage}
\end{tabular}
\end{tcolorbox}
}

%========================================================================================
% 끝
%========================================================================================


\begin{document}

\thispagestyle{firstpage}

\metainfo{Banking and Financial Institutions}{Module 1: Money and Finance}{Prof. Rustom Manouchehri Irani}{금융 시스템의 이해와 경제적 영향 분석}

\tableofcontents

\newpage

\section{개요}
\addcontentsline{toc}{section}{개요}
이 문서는 'Banking and Financial Institutions Module 1: Money and Finance'의 핵심 내용을 다룹니다.
돈의 정의와 기능, 지불 시스템의 진화, 통화량과 인플레이션의 관계를 탐구합니다.
또한 금융 시스템의 역할, 금융 시장과 금융 기관의 차이, 그리고 다양한 비은행 금융 기관의 기능을 이해합니다.
궁극적으로 금융 발전이 실물 경제 활동에 미치는 긍정적 및 부정적 영향을 종합적으로 파악하는 것이 목표입니다.
이 문서는 시험 대비를 위한 완결된 학습 자료로, 비전공자도 쉽게 이해할 수 있도록 상세한 설명과 예시를 제공합니다.

\newpage
\section{용어 정리}
\addcontentsline{toc}{section}{용어 정리}
다음은 본 모듈에서 다루는 주요 용어들을 비전공자도 이해하기 쉽게 정리한 표입니다.

\begin{adjustbox}{width=\textwidth,center}
\begin{tabular}{lllc}
\toprule
\textbf{용어 (Term)} & \textbf{쉬운 설명 (Easy Explanation)} & \textbf{원어 (Original)} & \textbf{비고 (Notes)} \\
\midrule
상품 화폐 & 그 자체로 가치가 있는 물건을 돈으로 사용하는 것 & Commodity Money & 조개껍데기, 소금, 금속 등 \\
법정 화폐 & 정부의 명령으로 가치가 부여된 돈 & Fiat Money & 오늘날 대부분의 지폐 \\
지불 시스템 & 돈을 주고받는 방법과 관련된 모든 제도와 인프라 & Payments System & 현금, 카드, 전자 결제 등 \\
M1 & 가장 좁은 의미의 통화량; 현금과 요구불 예금 & M1 (Money Supply 1) & 즉시 사용 가능한 돈 \\
M2 & M1에 저축성 예금 등을 더한 넓은 의미의 통화량 & M2 (Money Supply 2) & M1보다 유동성은 낮지만 쉽게 현금화 가능 \\
인플레이션 & 물가가 지속적으로 오르는 현상 & Inflation & 돈의 가치가 하락함을 의미 \\
금융 시스템 & 돈이 필요한 사람과 돈이 남는 사람을 연결하는 체계 & Financial System & 경제 성장의 핵심 인프라 \\
금융 상품 & 돈을 빌려주거나 투자할 때 주고받는 약속 증서 & Financial Instrument & 주식, 채권, 예금, 보험 등 \\
금융 시장 & 금융 상품이 거래되는 장소 & Financial Market & 주식 시장, 채권 시장 등 \\
금융 기관 & 돈을 모으고 빌려주는 역할을 하는 회사 & Financial Institution & 은행, 보험사, 연기금 등 \\
은행 & 예금을 받고 대출을 해주는 전통적인 금융 기관 & Bank & 상업 은행, 저축 기관, 신용 협동조합 \\
비은행 금융 기관 & 예금을 받지 않으면서 금융 서비스를 제공하는 기관 & Non-Bank Financial Institution & 보험사, 연기금, 증권사, 할부금융사 등 \\
유동성 & 자산을 현금으로 얼마나 쉽고 빠르게 바꿀 수 있는지 & Liquidity & 현금이 가장 유동성이 높음 \\
만기 불일치 & 돈을 빌려주는 기간과 빌리는 기간이 다른 문제 & Maturity Mismatch & 은행이 단기 예금으로 장기 대출을 해줌 \\
정보 비대칭 & 거래 당사자 중 한쪽이 더 많은 정보를 가진 상황 & Information Asymmetry & 대출자의 신용 정보를 은행이 더 잘 앎 \\
정부 후원 기업 & 정부가 설립했으나 민간 형태로 운영되는 기업 & Government Sponsored Enterprise (GSE) & Fannie Mae, Freddie Mac 등 \\
\bottomrule
\end{tabular}
\end{adjustbox}

\newpage
\section{핵심 개념 및 원리}
\addcontentsline{toc}{section}{핵심 개념 및 원리}

\subsection{1. 돈과 지불 시스템 (Money and Payments)}
\addcontentsline{toc}{subsection}{1. 돈과 지불 시스템}
이 섹션에서는 돈의 본질과 기능, 그리고 현대 사회에서 돈이 어떻게 교환되는지에 대한 지불 시스템을 탐구합니다. 돈은 우리가 매일 사용하는 개념이지만, 그 정의와 형태는 시대에 따라 크게 변화해 왔습니다.

\subsubsection*{1.1. 돈의 정의와 기능 (Money and Its Functions)}
\addcontentsline{toc}{subsubsection}{1.1. 돈의 정의와 기능}

\begin{tcolorbox}[summarybox]
  \textbf{핵심 요약:} 돈은 재화와 서비스의 지불 수단으로 널리 받아들여지는 특별한 자산이며, 교환의 매개, 가치 척도, 가치 저장의 세 가지 핵심 기능을 수행합니다.
\end{tcolorbox}

우리는 일상생활에서 항상 돈을 사용하지만, '돈이 무엇인가?'라는 질문은 생각보다 복잡합니다. 돈의 형태는 역사적으로 많이 변해왔고, 오늘날에도 그 정의에 대한 논의가 계속되고 있습니다.

\textbf{돈의 진화: 상품 화폐에서 법정 화폐까지}
돈의 역사를 살펴보면, 돈으로 사용된 자산들이 어떻게 변화해왔는지 알 수 있습니다.

\begin{itemize}
    \item \textbf{상품 화폐 (Commodity Money):}
    \begin{itemize}
        \item \textbf{한 줄 핵심 요약:} 그 자체로도 가치가 있는 물건을 돈으로 사용하는 형태입니다.
        \item \textbf{직관적 예시/비유:} 8세기 일본의 화살촉, 중국의 조개껍데기, 베네치아의 소금, 금과 같은 금속들이 상품 화폐의 예시입니다. 소금은 음식을 보존하는 데 유용하고, 화살촉은 사냥이나 방어에 사용될 수 있습니다.
        \item \textbf{기술적 정확한 설명:} 상품 화폐는 두 가지 특징을 가집니다. 첫째, 물리적인 물품에 의해 가치가 뒷받침됩니다. 둘째, 그 물품 자체에 본질적인 가치(intrinsic value)가 있습니다. 즉, 지불 수단으로 사용되지 않더라도 그 자체로 유용하거나 가치가 있는 것입니다. 16세기 스웨덴에서 사용된 구리판도 상품 화폐의 일종이었으나, 당시 구리의 가치가 낮아 거래를 위해 많은 양을 들고 다녀야 하는 불편함이 있었습니다.
    \end{itemize}

    \item \textbf{종이 화폐의 등장 (Paper Money):}
    \begin{itemize}
        \item \textbf{한 줄 핵심 요약:} 금속 화폐를 은행에 예치하고 그 증서로 거래하는 방식에서 시작되었습니다.
        \item \textbf{직관적 예시/비유:} 1657년 스웨덴의 Palmstruch가 설립한 Stockholms Banco는 유럽 최초로 종이 화폐를 발행했습니다. 이 종이 화폐는 은행에 예치된 구리 금속으로 언제든지 교환(redeemable on-demand)할 수 있었습니다. 사람들이 은행에 충분한 구리가 있다고 믿는 한, 종이 화폐는 가치를 가졌습니다.
        \item \textbf{기술적 정확한 설명:} 종이 화폐는 처음에는 예치된 상품 화폐(주로 금이나 은)에 의해 가치가 보증되었습니다. 미국도 1930년대까지 금본위제(gold standard)를 통해 종이 화폐를 금으로 교환할 수 있도록 했습니다. 그러나 Palmstruch 은행의 사례처럼, 은행이 충분한 금속을 보유하고 있지 않다는 소문이 돌면 사람들은 신뢰를 잃고 은행은 문을 닫을 수 있었습니다.
    \end{itemize}

    \item \textbf{법정 화폐 (Fiat Money):}
    \begin{itemize}
        \item \textbf{한 줄 핵심 요약:} 어떤 물리적 상품에도 의해 뒷받침되지 않고, 정부의 법적 강제력에 의해 가치가 부여되는 돈입니다.
        \item \textbf{직관적 예시/비유:} 오늘날 우리가 사용하는 지폐가 바로 법정 화폐입니다. 100달러 지폐를 인쇄하는 데 14센트밖에 들지 않지만, 우리는 그 지폐가 100달러의 가치를 가진다고 믿습니다. 이는 정부가 "이 지폐는 법정 통화이며, 모든 사적 사업체는 이를 지불 수단으로 받아들여야 한다"고 선언하고, 세금 납부에도 사용할 수 있기 때문입니다.
        \item \textbf{기술적 정확한 설명:} 법정 화폐는 그 자체로는 본질적인 가치가 없습니다. 그 가치는 전적으로 정부에 대한 신뢰와 법적 강제력에 기반합니다. 정부가 안정적이고 신뢰할 수 있는 한, 사람들은 법정 화폐를 가치 있는 지불 수단으로 받아들입니다.
    \end{itemize}
\end{itemize}

\textbf{현대 지불 수단의 진화}
돈의 형태는 종이 화폐를 넘어 플라스틱 카드와 디지털 방식으로 발전했습니다.

\begin{itemize}
    \item \textbf{직불 카드 (Debit Card):}
    \begin{itemize}
        \item \textbf{설명:} 내 은행 계좌에 있는 돈을 판매자에게 보내도록 은행에 지시하는 수단입니다. 계좌에 돈이 없으면 작동하지 않습니다.
    \end{itemize}
    \item \textbf{신용 카드 (Credit Card):}
    \begin{itemize}
        \item \textbf{설명:} 은행으로부터 대출을 받는 것입니다. 카드를 사용할 때 은행이 대신 돈을 지불하고, 나는 나중에 은행에 갚습니다. 1958년 Bank of America가 최초의 현대식 플라스틱 신용카드를 도입했습니다.
    \end{itemize}
    \item \textbf{스마트폰 앱 및 온라인 결제 시스템 (Smartphone Applications and Online Payment Systems):}
    \begin{itemize}
        \item \textbf{설명:} Apple Pay, Venmo, WeChat, PayPal 등은 플라스틱 카드 없이도 빠르고 안전하게 돈을 이체하거나 결제할 수 있게 해주는 21세기 시스템입니다.
    \end{itemize}
\end{itemize}

\textbf{돈의 공식적인 정의}
돈은 특정 자산으로, 재화와 서비스의 지불 또는 부채 상환(세금 납부 포함)을 위해 일반적으로 받아들여지는 수단입니다. 주머니 속 현금과 은행 계좌 잔액이 가장 중요한 예시입니다.

\begin{warningbox}
  \textbf{돈과 부(Wealth)의 차이:} 돈은 부와 같지 않습니다. 부유한 사람은 돈 외에도 집, 요트, 주식, 채권 등 다양한 자산을 가지고 있지만, 이들은 돈이 아닙니다. 돈은 오직 '지불 수단'으로서의 특별한 역할을 합니다.
\end{warningbox}

\textbf{돈의 세 가지 주요 기능}
돈이 지불 수단으로 기능하기 위해서는 다음과 같은 세 가지 특성을 가져야 합니다.

\begin{enumerate}
    \item \textbf{교환의 매개 (Medium of Exchange):}
    \begin{itemize}
        \item \textbf{한 줄 핵심 요약:} 거래를 원활하게 하고 물물교환의 불편함을 해소합니다.
        \item \textbf{직관적 예시/비유:} 내가 사과를 팔고 싶고 신발을 사고 싶을 때, 신발을 파는 사람이 사과를 원하지 않으면 거래가 어렵습니다 (물물교환의 이중 우연의 일치 문제). 돈이 있으면 사과를 팔아 돈을 받고, 그 돈으로 신발을 살 수 있습니다.
        \item \textbf{기술적 정확한 설명:} 현대 경제에서는 거래가 빈번하고 광범위하게 이루어집니다. 돈은 거래를 즉시 완료할 수 있게 하여, 지불을 기다릴 필요 없이 상품과 서비스를 교환할 수 있도록 합니다. 돈의 가치는 널리 이해되고 있어, 거래 시 상대방의 돈 가치를 일일이 확인할 필요가 없습니다.
    \end{itemize}

    \item \textbf{가치 척도 (Unit of Account):}
    \begin{itemize}
        \item \textbf{한 줄 핵심 요약:} 재화와 서비스의 가치를 측정하고 비교하는 기준이 됩니다.
        \item \textbf{직관적 예시/비유:} 모든 상품의 가격이 '달러와 센트'로 표시되어 있다면, 사과 한 개가 1달러이고 바나나 한 개가 50센트라는 것을 쉽게 비교할 수 있습니다.
        \item \textbf{기술적 정확한 설명:} 돈은 가격을 매기고 부채를 기록하는 단위가 됩니다. 이를 통해 다양한 제품의 가치를 쉽게 비교할 수 있으며, 경제 활동의 효율성을 높입니다.
    \end{itemize}

    \item \textbf{가치 저장 수단 (Store of Value):}
    \begin{itemize}
        \item \textbf{한 줄 핵심 요약:} 구매력을 현재에서 미래로 이전할 수 있게 해줍니다.
        \item \textbf{직관적 예시/비유:} 오늘 번 돈을 당장 쓰지 않고 저축했다가 나중에 휴가를 갈 때 사용할 수 있습니다. 돈은 그 가치를 비교적 안정적으로 유지해야 합니다.
        \item \textbf{기술적 정확한 설명:} 돈은 시간이 지나도 그 가치를 비교적 안정적으로 유지해야 합니다. 만약 돈의 가치가 불안정하게 급락한다면 (예: 초인플레이션), 상인들은 돈을 받으려 하지 않을 것입니다. 이 때문에 일부 국가에서는 자국 통화와 함께 미국 달러와 같은 안정적인 통화를 사용하기도 합니다.
    \end{itemize}
\end{enumerate}

\newpage
\subsubsection*{1.2. 지불 시스템 (The Payments System)}
\addcontentsline{toc}{subsubsection}{1.2. 지불 시스템}

\begin{tcolorbox}[summarybox]
  \textbf{핵심 요약:} 지불 시스템은 재화, 서비스, 자산의 교환을 가능하게 하는 인프라로, 현금 결제와 비현금 결제로 나뉘며, 비현금 결제는 전자화되는 추세입니다.
\end{tcolorbox}

돈과 지불 시스템은 밀접하게 연결되어 있습니다. 지불 시스템은 경제의 '금융 배관(financial plumbing)'이라고 불리며, 상품, 서비스, 자산의 교환을 가능하게 하는 일련의 체계입니다. 이 시스템은 우리가 인식하는 것보다 더 많은 인프라를 필요로 하며, 끊임없이 혁신되고 있습니다. 효율적인 지불 시스템은 경제에 필수적이므로 정부와 정책 입안자들은 이에 큰 관심을 기울입니다. 많은 경우, 연방준비제도(Federal Reserve System)와 같은 중앙은행은 지불을 촉진하기 위해 설립되었습니다.

\textbf{현금 결제와 비현금 결제}
모든 거래가 지폐 교환을 수반하는 것은 아닙니다. 우리는 비현금 결제(non-cash payments)를 많이 사용합니다. 모든 비현금 결제는 궁극적으로 돈의 교환을 수반하며, 주로 한 은행 계좌에서 다른 은행 계좌로 이루어집니다.

\textbf{주요 비현금 결제 유형}

\begin{enumerate}
    \item \textbf{수표 (Paper Checks):}
    \begin{itemize}
        \item \textbf{한 줄 핵심 요약:} 여전히 미국에서 중요한 지불 수단으로, 은행에 자금 이체를 지시하는 서면 명령입니다.
        \item \textbf{직관적 예시/비유:} 누군가에게 수표를 주면, 그 수표는 내 은행에 '내 계좌에서 판매자의 계좌로 돈을 옮겨라'는 지시를 내리는 것과 같습니다. 수표 자체는 돈이나 통화가 아닙니다.
        \item \textbf{기술적 정확한 설명:} 미국에서는 지역 연방준비은행이 종종 수표 청산(clearing)을 담당합니다. 이는 은행 계좌 간의 자금 이체를 처리하는 것을 의미합니다.
    \end{itemize}

    \item \textbf{자동 어음 교환 시스템 (Automated Clearing House, ACH) 거래:}
    \begin{itemize}
        \item \textbf{한 줄 핵심 요약:} 수표와 유사하게 은행 간 자금 이체를 지시하지만, 전자적으로 이루어져 더 빠르고 편리합니다.
        \item \textbf{직관적 예시/비유:} 직불카드나 신용카드로 결제할 때 주로 ACH 이체가 발생합니다. 직불카드는 내 계좌에서, 신용카드는 은행의 자금에서 이체가 이루어지며, 신용카드 사용 시에는 대출이 발생합니다.
        \item \textbf{기술적 정확한 설명:} ACH 이체는 전자적으로 처리되므로 수표보다 행정적 번거로움이 적고 빠릅니다. 선불 직불카드(prepaid debit cards)도 중요한 지불 수단이며, 일부 국가에서는 고객의 은행 계좌와 연결되지 않기도 합니다.
    \end{itemize}

    \item \textbf{디지털 지갑 및 송금 앱 (Digital Wallets and Money Transfer Apps):}
    \begin{itemize}
        \item \textbf{한 줄 핵심 요약:} 컴퓨터나 스마트폰을 통해 전자 직불/신용 이체를 시작할 수 있게 해주는 시스템입니다.
        \item \textbf{직관적 예시/비유:} Apple Pay는 시장 선두 주자이며, Venmo나 Zelle과 같은 앱은 친구나 가족 간의 송금을 가능하게 합니다. 이들은 주로 기존 직불/신용카드 정보를 저장하여 사용합니다.
        \item \textbf{기술적 정확한 설명:} 이러한 시스템은 결제 정보를 디지털 방식으로 저장하고, 사용자가 편리하게 전자 거래를 시작할 수 있도록 합니다.
    \end{itemize}
\end{enumerate}

\begin{tcolorbox}[summarybox]
  \textbf{핵심 요약:} 현금과 은행 계좌 잔액은 경제학자들이 돈의 핵심 척도로 간주합니다. 2018년 미국에서는 약 1,750억 건의 비현금 결제가 약 100조 달러 규모로 이루어졌으며, 전자 결제가 수표를 대체하는 추세입니다.
\end{tcolorbox}

\textbf{비현금 결제 방식별 가치 비중 (2012년 이후 추세)}
제공된 자료에 따르면, 비현금 결제 가치에서 수표가 차지하는 비중은 감소하고 전자 결제가 증가하는 추세입니다. 2012년에는 신용카드 결제액이 수표 결제액을 처음으로 넘어섰습니다. 그럼에도 불구하고 수표는 여전히 중요한 역할을 하고 있습니다.

\textbf{사례 연구: 비트코인은 돈인가? (Is Bitcoin Money?)}
암호화폐는 새로운 지불 시스템과 정부가 발행하지 않는 새로운 통화를 결합한 것입니다. 통화 자체는 사고팔 수 있는 디지털 토큰이며, 거래는 분산원장기술(distributed ledger technology), 즉 블록체인을 사용하여 기록됩니다. 블록체인은 정부 기관의 검증 없이 거래를 기록하고 검증하는 전자 데이터베이스입니다. 2008년부터 존재해 온 비트코인(Bitcoin)이 가장 유명한 예시입니다.

\textbf{비트코인의 잠재적 장점:}
\begin{itemize}
    \item \textbf{정부 통제 불가능:} 정부가 디지털 토큰의 공급을 통제하지 않으므로, 정부에 의해 가치가 훼손될 수 없습니다.
    \item \textbf{익명성:} 거래가 정부 기관의 검증을 필요로 하지 않으므로, 사용자는 익명으로 전자 결제를 할 수 있습니다.
\end{itemize}

\textbf{비트코인이 돈의 세 가지 특성을 충족하는가?}
\begin{enumerate}
    \item \textbf{유용한 지불 수단인가? (Means of Payment)}
    \begin{itemize}
        \item \textbf{평가:} 아직은 아닙니다. 상품 화폐처럼 본질적인 가치가 없고, 법정 화폐도 아니므로 기업이 비트코인을 의무적으로 받아들일 필요가 없습니다. 세금 납부에도 사용할 수 없습니다. 현재까지 비트코인 거래는 주로 마약 거래나 사이버 공격과 같은 범죄 활동과 연관되어 있습니다. 이러한 이유로 중국은 암호화폐 거래를 단속했습니다.
    \end{itemize}
    \item \textbf{안정적인 가치 저장 수단인가? (Stable Store of Value)}
    \begin{itemize}
        \item \textbf{평가:} 아닙니다. 비트코인의 가격은 크게 변동하여 통화라기보다는 투기성 자산에 가깝습니다. 미국 세법상 비트코인 및 기타 암호화폐는 '재산(property)'으로 분류되어 투자로 과세됩니다. 가치 변동이 심하기 때문에 택시 요금이나 테슬라 자동차 가격을 비트코인으로 책정한다면 매우 불안정할 것입니다.
    \end{itemize}
    \item \textbf{가치 척도인가? (Unit of Account)}
    \begin{itemize}
        \item \textbf{평가:} 위의 두 가지 이유로 인해, 비트코인은 안정적인 가치 척도로서의 역할을 제대로 수행하지 못합니다.
    \end{itemize}
\end{enumerate}
\begin{tcolorbox}[warningbox]
  \textbf{결론:} 현재로서는 비트코인이 미국 달러나 유로와 같은 신뢰할 수 있는 정부 발행 통화와는 거리가 멀다고 평가됩니다. 미래에는 바뀔 수도 있지만, 아직은 아닙니다.
\end{tcolorbox}

\newpage
\subsubsection*{1.3. 돈과 인플레이션 (Money and Inflation)}
\addcontentsline{toc}{subsubsection}{1.3. 돈과 인플레이션}

\begin{tcolorbox}[summarybox]
  \textbf{핵심 요약:} 통화량의 측정은 경제 성장에 영향을 미치므로 중요하며, 과거에는 통화량 증가가 인플레이션과 밀접한 관계를 보였으나, 21세기에는 그 관계가 불분명해졌습니다.
\end{tcolorbox}

이 세션에서는 돈의 개념을 넘어 경제에 유통되는 총 통화량을 측정하는 방법을 알아보고, 통화량과 인플레이션의 관계를 살펴봅니다. 통화량의 많고 적음은 경제 성장, 이자율, 인플레이션에 영향을 미칠 수 있으므로, 중앙은행과 정책 입안자들은 통화 공급과 지불 시스템의 혁신에 큰 관심을 기울입니다.

\textbf{돈과 유동성 (Money and Liquidity)}
돈은 우리가 지불에 사용하는 특별한 자산이며, 우리의 부(wealth)의 일부입니다. 우리는 거래를 위해 돈을 사용하지만, 우리의 부는 다양한 형태로 보유됩니다.

\begin{itemize}
    \item \textbf{자산 유동성 (Asset Liquidity):}
    \begin{itemize}
        \item \textbf{한 줄 핵심 요약:} 자산을 돈으로 전환하는 데 드는 비용과 시간을 나타냅니다.
        \item \textbf{직관적 예시/비유:} 현금은 이미 돈이므로 가장 유동성이 높습니다. 집이나 보석은 가치가 있지만, 적절한 가격에 팔려면 시간이 걸리므로 유동성이 낮습니다. 주식이나 저축 계좌는 그 중간 정도의 유동성을 가집니다.
        \item \textbf{기술적 정확한 설명:} 우리는 은행 계좌에 자금을 보유하고, 주식과 채권에 투자하며, 집을 소유할 수 있습니다. 이러한 다른 자산들도 지불에 사용될 수 있지만, 먼저 돈으로 전환(청산, liquidating)해야 하므로 더 많은 비용이 들 수 있습니다.
    \end{itemize}
\end{itemize}

\textbf{통화량 측정 (Measuring Money Supply)}
돈의 정의는 시간이 지남에 따라 변해왔으며, 이상적인 측정 방법은 없습니다. 현금만 돈으로 포함하는 것은 너무 엄격하고, 주식과 채권을 포함하는 것은 너무 느슨합니다. 경제학자들은 주로 두 가지 기본 정의인 M1과 M2를 사용합니다.

\begin{enumerate}
    \item \textbf{M1 (협의 통화):}
    \begin{itemize}
        \item \textbf{한 줄 핵심 요약:} 가장 엄격한 통화량 정의로, 대중이 보유한 현금과 은행의 요구불 예금(checking accounts)을 포함합니다.
        \item \textbf{직관적 예시/비유:} 내 지갑에 있는 지폐와 언제든지 인출하거나 직불카드로 사용할 수 있는 은행 계좌 잔액이 M1에 해당합니다.
        \item \textbf{기술적 정확한 설명:} M1은 즉시 지불 수단으로 사용될 수 있는 가장 유동적인 자산들로 구성됩니다.
    \end{itemize}

    \item \textbf{M2 (광의 통화):}
    \begin{itemize}
        \item \textbf{한 줄 핵심 요약:} M1에 일부 저축성 예금(savings accounts)을 추가한 더 넓은 통화량 측정치입니다.
        \item \textbf{직관적 예시/비유:} M1에 더해, 은행의 저축 계좌나 머니마켓 펀드(money market fund)와 같이 M1보다는 유동성이 약간 낮지만 쉽게 현금화할 수 있는 자금들이 M2에 포함됩니다. 저축 계좌는 이자를 지급하지만, 월별 인출 횟수 제한이 있을 수 있습니다.
        \item \textbf{기술적 정확한 설명:} 가계가 이자를 지급하는 저축 계좌로 자금을 옮기면서 M2가 경제 활동과 더 밀접하게 움직이는 경향이 있어, 경제학자들은 M2에 더 주목하는 경향이 있습니다.
    \end{itemize}
\end{enumerate}

\begin{tcolorbox}[examplebox]
  \textbf{통화량 구성 예시 (2020년 4월 미국 연방준비제도 데이터):}
  \begin{itemize}
      \item M1은 통화(currency)와 요구불 예금(checking accounts)으로 구성됩니다.
      \item M2는 M1에 저축성 예금(savings accounts)을 더한 것입니다.
  \end{itemize}
  만약 현금만 돈으로 측정한다면 전체 그림의 큰 부분을 놓치게 됩니다. 그래서 경제학자들은 M2를 선호합니다.
\end{tcolorbox}

\textbf{인플레이션 (Inflation)}
인플레이션은 시간이 지남에 따라 물가가 상승하는 비율을 나타냅니다. 인플레이션이 높으면 같은 상품을 사기 위해 더 많은 돈을 지불해야 합니다. 즉, 돈의 가치가 하락합니다.

\begin{itemize}
    \item \textbf{소비자 물가 지수 (Consumer Price Index, CPI):}
    \begin{itemize}
        \item \textbf{한 줄 핵심 요약:} 인플레이션율을 측정하는 가장 일반적인 지표로, 특정 기간 동안 동일한 재화 및 서비스 바구니(식료품, 주택, 휴대폰 등) 가격의 변화를 측정합니다.
        \item \textbf{직관적 예시/비유:} 2021년 1월 미국 CPI가 1.4\%였다면, 2020년 1월에 100달러였던 상품 바구니가 2021년 1월에는 101.40달러가 되었다는 의미입니다.
        \item \textbf{기술적 정확한 설명:} CPI의 성장은 돈의 구매력에 대한 정보를 제공합니다.
    \end{itemize}
\end{itemize}

\textbf{초인플레이션 (Hyperinflation)}
과거에는 매우 높은 인플레이션(초인플레이션) 사례들이 있었습니다.
\begin{itemize}
    \item 1923년 독일: 일일 인플레이션율 21\%, 3일마다 물가 두 배 상승.
    \item 1946년 헝가리, 2008년 짐바브웨, 2018년 베네수엘라에서도 유사한 사례 발생.
\end{itemize}
\begin{tcolorbox}[warningbox]
  \textbf{초인플레이션의 원인:} 일반적으로 정부가 재정 적자를 메우기 위해 중앙은행에 돈을 인쇄하도록 강요하여 통화 공급을 급격히 늘릴 때 발생합니다. 고정된 수의 상품을 쫓는 돈이 점점 많아지면 돈의 가치는 하락하고, 이는 가계와 기업에 심각한 결과를 초래합니다.
\end{tcolorbox}

\textbf{통화량 증가와 인플레이션의 관계}

\begin{enumerate}
    \item \textbf{1960년~1980년: 강한 양의 상관관계}
    \begin{itemize}
        \item \textbf{설명:} 1970년대와 80년대의 '대인플레이션(The Great Inflation)' 시기에는 인플레이션율이 10\%를 초과했습니다. 이 시기(1960~1980년)의 데이터 분석 결과, M2 통화량 증가율(2년 시차 적용)과 CPI 인플레이션율 사이에 강한 양의 상관관계가 나타났습니다. 즉, 높은 통화량 증가는 2년 후 높은 인플레이션과 연관되었고, 낮은 통화량 증가는 낮은 인플레이션과 연관되었습니다. 이는 우리의 직관과 일치합니다.
    \end{itemize}
    \item \textbf{2000년 이후: 상관관계 약화}
    \begin{itemize}
        \item \textbf{설명:} 2000년 이후 기간에 초점을 맞추면, M2 증가율과 CPI 인플레이션율 사이의 상관관계가 사라집니다. M2는 21세기 인플레이션을 예측하는 데 더 이상 유용하지 않습니다.
        \item \textbf{원인 가설:}
        \begin{enumerate}
            \item \textbf{낮은 인플레이션 환경:} 2000년 이후 인플레이션율이 역사적 기준으로 비정상적으로 낮은 약 2\% 수준을 유지했기 때문에, 통화량 증가와 인플레이션 간의 관계가 명확하게 나타나지 않을 수 있습니다.
            \item \textbf{돈의 측정 방식 변화:} 21세기에는 금융 시장의 거래 비용이 낮아지고, 신용카드가 널리 보급되었으며, 새로운 디지털 결제 기술이 매일 등장하면서 돈을 사용하고 지불하는 방식이 크게 변했습니다. 따라서 M2가 더 이상 돈을 정확하게 측정하는 지표가 아닐 수 있습니다. 새로운 통화량 측정 방식이 필요할 수도 있습니다.
        \end{enumerate}
    \end{itemize}
\end{enumerate}

\begin{tcolorbox}[summarybox]
  \textbf{결론:} 돈의 개념에서 통화량 측정으로 나아가, 통화량 증가와 인플레이션의 관계를 살펴보았습니다. 1970년대와 80년대의 대인플레이션 시기에는 통화량 증가가 높았지만, 21세기에는 이 관계가 불분명해졌습니다. 이는 미래에 정부와 중앙은행이 인플레이션을 통제하는 과거의 접근 방식이 효과적일지에 대한 의문을 제기합니다.
\end{tcolorbox}

\newpage
\subsection{2. 금융 시스템 (The Financial System)}
\addcontentsline{toc}{subsection}{2. 금융 시스템}
이 섹션에서는 금융 시스템이 경제에서 어떤 역할을 하는지, 그리고 금융 시장과 금융 기관이 어떻게 자금을 중개하는지 상세히 살펴봅니다.

\subsubsection*{2.1. 금융 시스템의 역할 (What Does the Financial System Do?)}
\addcontentsline{toc}{subsubsection}{2.1. 금융 시스템의 역할}

\begin{tcolorbox}[summarybox]
  \textbf{핵심 요약:} 금융 시스템은 자금 여유가 있는 '저축자'와 자금 수요가 있는 '차입자'를 연결하여 자금을 가장 효율적인 곳으로 흐르게 하고, 그 대가로 '금융 상품'을 교환함으로써 경제 번영과 성장을 촉진합니다.
\end{tcolorbox}

금융 시스템의 출발점은 자금 수요가 있는 경제 주체입니다. 이는 수익성 있는 투자 기회를 가진 기업일 수도 있고, 주택 구매 자금이 필요한 가계일 수도 있습니다.

\textbf{저축자와 차입자 (Savers and Borrowers)}

\begin{itemize}
    \item \textbf{차입자 (Borrowers):}
    \begin{itemize}
        \item \textbf{한 줄 핵심 요약:} 자금이 필요한 개인, 기업, 정부를 포함합니다.
        \item \textbf{직관적 예시/비유:} 주택을 사거나 물건을 구매하기 위한 가계, 단기 운전자금이나 장기 사업 확장을 위한 기업, 공공재를 위한 정부, 트랙터를 사려는 농부 등이 차입자가 될 수 있습니다.
        \item \textbf{기술적 정확한 설명:} 기업은 자체 현금 보유액이나 자산 매각으로 투자 자금을 조달할 수 있지만, 내부 자금이 충분하지 않을 경우 외부에서 자금을 조달해야 합니다. 이때 저축자로부터 자금을 빌리는 차입자가 됩니다.
    \end{itemize}

    \item \textbf{저축자 (Savers):}
    \begin{itemize}
        \item \textbf{한 줄 핵심 요약:} 현재 지출보다 더 많은 소득을 가진 가계가 대부분입니다.
        \item \textbf{직관적 예시/비유:} 은퇴를 위해 연기금에 돈을 넣거나, 은행에 저축하거나, 비상금을 위해 보험에 가입하는 가계가 저축자입니다.
        \item \textbf{기술적 정확한 설명:} 저축은 미래를 위한 준비이며, 이는 연기금이나 뮤추얼 펀드를 통해 이루어지거나, 은행 예금으로 현금 관리에 활용되거나, 보험 가입을 통해 위험을 분산하는 형태로 나타납니다.
    \end{itemize}
\end{itemize}

\textbf{금융 시스템의 역할}
금융 시스템은 저축자와 차입자를 연결하여 저축된 자금을 차입자의 자금 수요에 맞게 전달합니다. 이 과정의 궁극적인 목표는 경제적 번영입니다. 금융 시스템은 자원을 가장 효율적인 용도로 배분하여 경제 성장을 이끌어냅니다. 저축자들은 자금을 제공하는 대가로 '금융 상품(financial instruments)'을 받습니다.

\textbf{금융 상품 (Financial Instruments)}
금융 상품은 기본적으로 '상환에 대한 공식적인 약속'입니다. 이는 법적 구속력이 있는 의무이며, 문제가 발생할 경우 법원과 정부가 강제할 수 있습니다. 특정 상황에서 참여자들 간에 자금이 이전되도록 요구합니다.

\begin{itemize}
    \item \textbf{저축자를 위한 금융 상품:}
    \begin{itemize}
        \item \textbf{은행 예금 계좌 (Bank Deposit Account):}
        \begin{itemize}
            \item \textbf{설명:} 저축자는 은행에 현금을 예치하고, 이 자금은 요구불(on-demand)로 인출 가능합니다. 안전하고 편리하며, 수표, 직불카드, 모바일 뱅킹을 통해 결제에 사용될 수 있습니다.
        \end{itemize}
        \item \textbf{뮤추얼 펀드 주식 (Mutual Fund Share):}
        \begin{itemize}
            \item \textbf{설명:} 돈을 투자 포트폴리오의 지분과 교환하여 부를 저장하고 성장시키는 데 도움을 줍니다.
        \end{itemize}
        \item \textbf{보험 (Insurance):}
        \begin{itemize}
            \item \textbf{설명:} 사고와 같은 불리한 사건에 대비하여 현재 지출을 포기하고 보험료를 지불합니다. 사고 발생 시 보험금(payout)을 받아 손실을 보전합니다. 보험은 저축자가 위험을 이전할 수 있게 해주는 금융 상품입니다.
        \end{itemize}
    \end{itemize}

    \item \textbf{차입자를 위한 금융 상품:}
    \begin{itemize}
        \item \textbf{주식 (Equity):}
        \begin{itemize}
            \item \textbf{설명:} 기업 소유주가 주식을 팔아 자금을 조달합니다. 주식 구매자는 기업의 지분을 소유하게 되며, 이는 미래 기업 이익에 대한 권리를 부여합니다.
        \end{itemize}
        \item \textbf{채무 (Debt):}
        \begin{itemize}
            \item \textbf{설명:} 채권(bonds), 은행 대출(bank loans), 모기지(mortgages) 등이 있습니다. 차입자는 오늘 일시불로 자금을 받고, 미래에 정해진 날짜에 정해진 금액을 상환하기로 약속합니다. 이 때문에 '고정 수입(fixed income)' 상품이라고도 불립니다.
            \item \textbf{예시:} 정부 채권은 과거에 쿠폰이 붙어 있어 은행에서 현금으로 교환할 수 있었습니다. 주택 담보 대출(residential mortgage)의 경우, 차입자는 토지와 주택을 담보(collateral)로 제공합니다.
        \end{itemize}
    \end{itemize}
\end{itemize}

\begin{tcolorbox}[warningbox]
  \textbf{왜 이렇게 많은 금융 상품이 존재할까?}
  다양한 금융 상품은 저축자와 차입자의 다양한 요구를 충족시키기 위해 맞춤 제작됩니다. 저축자는 편리한 부의 저장 및 성장 방법을 원하고, 위험을 분산하기 위해 보험을 원할 수 있습니다. 차입자는 장기 자금을 원하며, 부채 발행이나 주식 매각을 통해 자금을 조달할 수 있습니다.
\end{tcolorbox}

\begin{tcolorbox}[summarybox]
  \textbf{결론:} 금융 시스템은 자금 여유가 있는 저축자와 투자 수요가 있는 차입자 사이에 존재합니다. 이 시스템은 저축자로부터 차입자에게 자금을 전달하고, 그 대가로 저축자와 차입자의 요구를 충족시키도록 설계된 금융 상품을 교환합니다.
\end{tcolorbox}

\newpage
\subsubsection*{2.2. 금융 시장 vs. 금융 기관 (Financial Markets Versus Financial Institutions)}
\addcontentsline{toc}{subsubsection}{2.2. 금융 시장 vs. 금융 기관}

\begin{tcolorbox}[summarybox]
  \textbf{핵심 요약:} 금융 시스템은 자금을 저축자에서 차입자로 전달하는 두 가지 주요 경로를 가집니다. 하나는 '금융 시장'을 통한 직접적인 경로이고, 다른 하나는 '금융 기관'을 통한 간접적인 경로입니다.
\end{tcolorbox}

금융 시스템은 저축자로부터 차입자에게 자금을 전달하고, 그 대가로 금융 상품을 교환합니다. 이 과정은 크게 두 가지 방식으로 이루어집니다.

\textbf{1. 금융 시장을 통한 직접 자금 흐름 (Direct Flow via Financial Markets)}
저축자가 금융 시장에 직접 참여하는 방식입니다.

\begin{itemize}
    \item \textbf{한 줄 핵심 요약:} 저축자가 주식이나 채권과 같은 금융 상품을 다른 투자자로부터 직접 구매하는 방식입니다.
    \item \textbf{직관적 예시/비유:} 내가 Apple 주식을 사고 싶다면, 보통 Apple 회사로부터 직접 사는 것이 아니라 뉴욕 증권 거래소(New York Stock Exchange)에서 다른 투자자로부터 주식을 구매합니다.
    \item \textbf{기술적 정확한 설명:} 이러한 '직접 교환(direct exchange)'은 보통 브로커(broker)나 딜러(dealer)와 같은 금융 회사의 도움을 받아 이루어집니다.
    \begin{itemize}
        \item \textbf{브로커 (Broker):} 수수료를 받고 거래 파트너를 찾아줍니다.
        \item \textbf{딜러 (Dealer):} 수수료를 받고 자신의 재고(inventory)에 있는 금융 상품을 직접 판매하여 거래 파트너 역할을 합니다.
    \end{itemize}
    오늘날 대부분의 이러한 상호작용은 E*Trade와 같은 온라인 거래 플랫폼을 통해 전자적으로 이루어집니다.
\end{itemize}

\textbf{금융 시장의 종류}
금융 시장은 거래되는 금융 상품의 종류에 따라 분류됩니다.

\begin{itemize}
    \item \textbf{주식 시장 (Stock or Equity Markets):} 기업의 주식이 거래되는 곳입니다.
    \item \textbf{채무 시장 (Debt Markets):} 대출, 모기지, 회사채 등이 거래되는 곳입니다.
    \item \textbf{단기 금융 시장 (Money Markets):} 만기가 1년 미만인 단기 채권(예: 기업어음, commercial paper)이 거래되는 곳입니다.
\end{itemize}
\begin{tcolorbox}[examplebox]
  \textbf{Apple의 자금 조달 예시:} Apple이 운영 자금을 조달할 때, 주식, 장기 채권, 또는 단기 채권(기업어음)을 발행할 수 있습니다. 발행 후, 주식은 주식 시장에서, 장기 채권은 채권 시장에서, 기업어음은 단기 금융 시장에서 거래됩니다.
\end{tcolorbox}

\textbf{잘 작동하는 금융 시장의 특징}
잘 작동하는 금융 시장은 현대 경제에서 핵심적인 역할을 합니다. 크고 활발한 시장을 가진 국가들은 다음 세 가지 특징을 보입니다.

\begin{enumerate}
    \item \textbf{낮은 거래 비용 (Low Transaction Costs):} 차입자가 증권을 발행하는 비용과 투자자들이 서로 거래하는 비용이 저렴해야 합니다.
    \item \textbf{정보의 효율적 흐름 (Efficient Information Flow):} 시장 참여자들이 차입자에 대한 정보를 수집, 처리, 전달할 수 있어야 합니다. 이는 시장 가격이 잠재적 위험과 수익을 정확하게 반영하도록 보장하며, 자원이 최적의 용도로 흐르도록 합니다.
    \item \textbf{투자자 보호 (Investor Protection):} 투자자들이 투자 손실로부터 보호받는다는 의미가 아니라, 차입자가 투자자를 속이지 않을 것이라는 신뢰를 가져야 한다는 의미입니다. 금융 상품은 법적 구속력이 있는 상환 약속이므로, 현대 금융 시장은 규제 기관과 법원이 계약을 강제할 것을 요구합니다. 투자자들이 남용으로부터 보호받을 때 기꺼이 투자할 것입니다.
\end{enumerate}

\textbf{2. 금융 기관을 통한 간접 자금 흐름 (Indirect Flow via Financial Institutions)}
금융 기관은 저축자로부터 자금을 조달하고, 그 자금을 차입자에게 투자하는 방식으로 자금을 중개합니다.

\begin{itemize}
    \item \textbf{한 줄 핵심 요약:} 금융 기관은 저축자로부터 자금을 모아 단기 부채(예: 예금)를 제공하고, 이 자금을 차입자에게 장기 투자(예: 대출)하는 방식으로 자금을 중개합니다.
    \item \textbf{기술적 정확한 설명:} 금융 기관은 먼저 저축자로부터 자금을 조달하고, 그 대가로 보통 단기 부채를 제공합니다. 그런 다음, 이 자금을 차입자가 발행한 금융 상품에 투자합니다. 이 과정에서 저축자로부터 차입자로 자금이 전달되지만, 중간에 금융 기관이라는 추가 단계가 있습니다. 브로커나 딜러와 달리, 금융 기관은 대차대조표(balance sheets)를 가지고 있으며 종종 장기 투자를 합니다.
\end{itemize}

\textbf{은행의 예시를 통한 금융 기관의 역할}
가장 단순화된 은행의 예를 통해 이 과정을 이해해 봅시다.

\begin{enumerate}
    \item \textbf{저축자의 예금:} 내가 저축자로서 1,000달러를 은행에 예치합니다. 은행은 나에게 '예금 계약(deposit contract)'을 제공하며, 이는 은행의 부채(liability)가 됩니다 (은행은 내가 원할 때 1,000달러를 돌려줄 의무가 있습니다). 은행은 이러한 예금을 여러 번 반복하여 총 100만 달러의 저축을 모읍니다.
    \item \textbf{은행의 대출:} 은행은 신용도 높은 사업체를 찾아 100만 달러를 장기 사업 대출로 빌려줍니다. 이 사업 대출은 은행의 자산(asset)이 됩니다 (사업체는 은행에 100만 달러를 갚을 의무가 있습니다).
    \item \textbf{자금 중개 및 대차대조표 형성:} 결과적으로 은행은 저축자로부터 차입자에게 자금을 전달했습니다. 이 과정에서 은행은 100만 달러의 대출(자산)과 100만 달러의 예금(부채)으로 구성된 대차대조표를 생성합니다.
\end{enumerate}

\textbf{왜 금융 기관이 필요한가? (금융 기관의 가치 창출)}
금융 시장과 은행이 공존하는 것을 보면, 둘 다 존재할 만한 좋은 이유가 있을 것입니다. 경제학자들은 금융 기관이 다음 세 가지 문제를 해결하는 데 도움을 준다고 봅니다.

\begin{enumerate}
    \item \textbf{만기 불일치 (Maturity Mismatch):}
    \begin{itemize}
        \item \textbf{한 줄 핵심 요약:} 차입자는 장기 자금을 원하고, 저축자는 단기 유동성을 원하는데, 은행이 이 둘을 연결해줍니다.
        \item \textbf{직관적 예시/비유:} 주택 구매자는 15년 또는 30년 만기의 모기지를 원하고, 기업은 장기 회수 기간을 가진 프로젝트를 위한 자금을 원합니다. 반면 저축자들은 자신의 자금에 쉽게 접근하고 싶어 하며, 모든 돈이 장기간 묶이는 것을 원치 않습니다.
        \item \textbf{기술적 정확한 설명:} 은행은 요구불 예금과 같은 단기 부채를 발행하고, 이를 장기 대출로 전환함으로써 이러한 만기 불일치 문제를 해결합니다.
    \end{itemize}

    \item \textbf{가격 위험 (Price Risk):}
    \begin{itemize}
        \item \textbf{한 줄 핵심 요약:} 주식이나 채권 가격의 변동성으로 인한 위험을 은행 예금은 줄여줍니다.
        \item \textbf{직관적 예시/비유:} 주식과 채권의 가격은 매일 변동하며, 거래 수수료도 변동할 수 있습니다. 일부 저축자들은 이러한 가격 위험을 매우 싫어하여 현금을 보유하는 것을 선호할 수 있습니다.
        \item \textbf{기술적 정확한 설명:} 은행의 예금 계약은 매우 편리하고 안전합니다. 정부가 제공하는 예금 보험이 있는 경우, 오늘 은행에 1달러를 예치하면 내일 최소 1달러를 인출할 수 있습니다. 은행 계좌와 관련된 수수료는 있을 수 있지만, 가격 위험은 없습니다.
    \end{itemize}

    \item \textbf{정보 비대칭 (Information Asymmetry):}
    \begin{itemize}
        \item \textbf{한 줄 핵심 요약:} 개별 저축자가 투자 위험과 수익을 평가하기 어려운 문제를 금융 기관이 전문적으로 해결해줍니다.
        \item \textbf{직관적 예시/비유:} Zano라는 아기 드론 프로젝트는 2015년 크라우드펀딩으로 200만 달러 이상을 모았지만, 결국 실패하고 투자자들은 돈을 잃었습니다. 개별 투자자들은 성공적인 사업을 선별하고 손실을 최소화하기 어렵습니다.
        \item \textbf{기술적 정확한 설명:} 정보 비대칭은 저축자가 투자에 대한 위험과 기대 수익을 평가하기 어려울 때 발생합니다. 은행은 차입자를 심사하고 평가하는 전문가이며, 대출 후에도 차입자를 모니터링하고 실적이 좋지 않은 차입자에게는 대출을 중단함으로써 이러한 정보 위험을 해결합니다. 시장 기반 솔루션(예: 주식 분석가, 신용 평가 기관)도 있지만, 정보 위험이 너무 크면 일부 가계는 투자를 꺼릴 수 있습니다.
    \end{itemize}
\end{enumerate}

\begin{tcolorbox}[summarybox]
  \textbf{결론:} 금융 시스템은 저축자로부터 차입자에게 자금을 이동시키며, 이는 두 가지 경로를 통해 이루어집니다. 첫째, 금융 시장을 통한 직접적인 경로로, 브로커의 도움을 받아 투자자들이 서로 증권을 거래합니다. 둘째, 금융 기관을 통한 간접적인 경로로, 금융 기관은 만기 불일치, 가격 위험, 정보 비대칭 문제를 해결하여 가치를 더합니다.
\end{tcolorbox}

\newpage
\subsubsection*{2.3. 은행 vs. 비은행 금융 기관 (Banks Versus Non-Banks)}
\addcontentsline{toc}{subsubsection}{2.3. 은행 vs. 비은행 금융 기관}

\begin{tcolorbox}[summarybox]
  \textbf{핵심 요약:} 현대 금융 산업은 금융 시장과 금융 기관으로 구성되며, 금융 기관은 다시 '은행'과 '비은행 금융 기관'으로 나뉩니다. 은행은 예금을 받고 대출을 하는 전통적인 역할을 하며, 비은행 기관은 예금을 받지 않으면서 다양한 금융 서비스를 제공합니다.
\end{tcolorbox}

이 세션에서는 미국 금융 산업의 조직 방식과 가장 중요한 은행 및 비은행 금융 기관들을 소개하고, 경제에서 이들의 역할을 설명합니다.

\textbf{금융 산업의 구성}
현대 금융 산업은 크게 금융 시장과 금융 기관으로 구성됩니다.
\begin{itemize}
    \item \textbf{금융 시장:} 저축자와 차입자를 직접 연결하는 중개 서비스를 제공하는 회사들이 있습니다.
    \item \textbf{금융 기관:} 저축자로부터 자금을 조달하여 주식, 대출, 채권, 부동산 등 다양한 금융 상품에 투자하는 회사들입니다. 이들은 대차대조표를 가집니다.
\end{itemize}
우리의 목표는 가장 중요한 금융 기관들이 누구이며, 어떤 고유한 경제적 기능을 수행하고, 어떻게 자금을 조달하며, 어디에 자금을 배치하는지(즉, 대차대조표가 어떻게 생겼는지) 이해하는 것입니다. 또한 시장 점유율 측면에서 어떤 기관이 가장 중요한지도 살펴봅니다.

\textbf{은행과 비은행 금융 기관의 구분}
금융 기관은 크게 '은행(banks)'과 '비은행 금융 기관(non-banks)'으로 나뉩니다.

\textbf{1. 은행 (Banks)}
\begin{itemize}
    \item \textbf{한 줄 핵심 요약:} 우리가 일상적으로 접하는 금융 기관으로, 저축자로부터 예금을 받고 차입자에게 대출을 해주는 전통적인 역할을 합니다.
    \item \textbf{종류:} 상업 은행(commercial banks), 저축 기관(savings institutions), 신용 협동조합(credit unions) 등이 있습니다. 다른 나라에는 다른 유형의 은행도 있지만, 모두 예금을 모으고 대출을 해주는 동일한 기능을 수행합니다.
    \item \textbf{자금 조달 (Funding):} 주로 개인 고객(저축자)으로부터 예금 계좌 형태로 자금을 조달합니다. 일부 대형 은행은 은행 간 시장(interbank market)이나 환매 조건부 채권 시장(repo market)에서 차입하기도 합니다.
    \item \textbf{자금 운용 (Deployment):} 주로 대출에 투자합니다. 사실 은행은 경제에서 대출의 주요 원천입니다. 은행은 어떤 가계가 모기지를 받을지, 대출 규모는 얼마인지, 이자율은 얼마인지를 결정합니다.
    \item \textbf{은행의 특별한 점 (가치 창출):}
    \begin{itemize}
        \item \textbf{예금자에게:} 안전하고 편리하며 지불 시스템에 접근할 수 있는 은행 계좌를 제공합니다.
        \item \textbf{은행 자체:} 소액의 저축을 대량으로 모아 신용도 높은 차입자에게 대출을 해줍니다. 이러한 '대출 원천화(origination)'는 은행의 핵심 기능입니다.
        \item \textbf{고객을 위한 시장 거래:} 고객을 대신하여 금융 시장에서 거래를 수행하여 위험을 분산하고 부를 관리할 수 있도록 돕습니다.
    \end{itemize}
\end{itemize}

\textbf{은행과 증권 시장의 상대적 중요성 (전 세계 데이터)}
세계은행(World Bank)의 최근 데이터에 따르면, 국내 증권 시장(주식 시장 시가총액 + 발행 채권 총액)의 규모와 총 대출 잔액을 비교하여 은행의 상대적 중요성을 파늠할 수 있습니다.

\begin{itemize}
    \item \textbf{주요 시사점 (지난 25년간 미국, 유럽, 아시아 두 경제권 데이터):}
    \begin{enumerate}
        \item \textbf{모두 중요:} 모든 경제권에서 증권 시장과 대출 시장 모두 중요합니다. 이는 은행이 여전히 매우 중요하다는 것을 의미합니다.
        \item \textbf{지역별 차이:} 전 세계적으로 많은 차이가 있습니다. 미국은 금융이 은행에 덜 의존하는 반면, 중국은 그 반대입니다.
    \end{enumerate}
\end{itemize}

\textbf{2. 비은행 금융 기관 (Non-Banks)}
\begin{itemize}
    \item \textbf{한 줄 핵심 요약:} 정부 보증 예금을 받지 않는 금융 기관입니다.
    \item \textbf{설명:} 비은행 금융 기관은 일반 은행과 유사한 많은 경제적 기능을 수행합니다. 부를 저장하는 안전하고 편리한 방법을 제공하고, 소액 저축을 모아 차입자에게 전달하며, 투자와 관련된 정보 비용을 최소화하고, 고객이 위험을 줄일 수 있도록 돕습니다.
    \item \textbf{종류:} 보험 회사(insurance companies), 연기금(pension funds), 증권 회사(securities firms), 할부금융 회사(finance companies), 정부 후원 기업(government-sponsored enterprises, GSEs) 등이 있습니다. 이들은 매우 다양하며, 자금 조달 및 투자 프로필, 그리고 존재 이유가 크게 다릅니다.
\end{itemize}

\begin{tcolorbox}[warningbox]
  \textbf{비은행 금융 기관의 광범위함:} P2P 대출업체부터 전당포, 사모펀드, 그리고 최근의 핀테크 기업까지 매우 광범위한 비은행 금융 기관 생태계가 존재합니다.
\end{tcolorbox}

\newpage
\subsubsection*{2.4. 비은행 금융 기관 개요 (Overview of Non-Banks)}
\addcontentsline{toc}{subsubsection}{2.4. 비은행 금융 기관 개요}

이 세션에서는 미국에서 자산 관리 규모 기준으로 가장 중요한 비은행 금융 기관들에 초점을 맞춥니다.

\textbf{1. 보험 회사 (Insurance Companies)}
\begin{itemize}
    \item \textbf{경제적 기능:} 위험 완화(risk mitigation)입니다. 개인의 관점에서는 '비상금을 위한 저축'의 한 형태입니다.
    \item \textbf{작동 방식:} 나는 어떤 불리한 사건의 위험을 보험 회사에 이전합니다. 보험료(premium)를 지불하고, 나쁜 상황(예: 자동차 사고)이 발생하면 보험금을 청구하여 보험 회사가 손실을 보전해줍니다. 보험사 입장에서는 더 많은 보험 계약을 체결할수록 보험금 지급이 예측 가능해집니다.
    \item \textbf{주요 유형:}
    \begin{itemize}
        \item \textbf{생명 보험 (Life Insurance):} 장애나 사망으로 인한 소득 손실로부터 보호합니다. (예: 19세기 나폴레옹 전쟁 참전 군인 가족을 위한 Scottish Widows)
        \item \textbf{손해 보험 (Property and Casualty Insurance):} 홍수, 화재와 같은 자연재해나 사고로 인한 가계 및 기업의 손실로부터 보호합니다. (예: 자동차 사고)
    \end{itemize}
    \item \textbf{자금 조달 및 운용:} 보험료는 모아져 투자됩니다. 보험 회사의 자금원은 모아진 보험료이며, 이 자금은 보험금 지급에 사용되지만, 주로 주식, 부동산, 채권과 같은 장기 투자에도 사용됩니다. 보험 회사는 미국 회사채의 가장 큰 투자자 중 하나입니다.
\end{itemize}

\textbf{2. 연기금 (Pension Funds)}
\begin{itemize}
    \item \textbf{경제적 기능:} 미래를 위해 부를 저장하고 성장시키는 데 사용됩니다.
    \item \textbf{작동 방식:} 가계는 젊은 나이부터 은퇴를 위해 저축할 수 있으며, 종종 강제적으로 이루어지기도 합니다 (예: 급여의 일부가 자동 적립). 은퇴 저축 계좌는 일반적으로 소득세 유예와 같은 세금 혜택을 제공하며, 고용주가 기여금을 내기도 합니다. 은퇴 시에는 일시불 또는 종신 연금 형태로 지급받습니다.
    \item \textbf{자금 조달 및 운용:} 연기금은 직원 및 고용주의 기여금을 모아 자금을 조달합니다. 이 자금은 주로 채권, 주식, 부동산과 같은 장기 투자에 사용되며, 동시에 은퇴자에게 지급할 보험금을 관리합니다.
    \item \textbf{보험 회사와의 유사점:} 연기금과 생명 보험 회사는 매우 유사한 사업을 합니다. 연기금은 은퇴할 경우 지급하고, 생명 보험은 사망할 경우 지급합니다. 이러한 유사성 때문에 이들 사업은 동일한 투자를 하는 경향이 있으며, 대형 생명 보험 회사가 연기금을 관리하기도 합니다.
\end{itemize}

\textbf{3. 증권 회사 (Securities Firms)}
주요 유형은 투자 은행과 투자 펀드입니다.

\begin{itemize}
    \item \textbf{투자 은행 (Investment Banks):}
    \begin{itemize}
        \item \textbf{한 줄 핵심 요약:} 상업 은행이 가계의 대출을 결정하는 것과 달리, 투자 은행은 기업이 자본 시장에서 증권을 발행할 수 있는지와 그 조건을 결정합니다.
        \item \textbf{핵심 기능: 인수 (Underwriting):} 기업이 증권을 발행하고자 할 때, 투자 은행은 특정 가격으로 전체 발행 물량을 매입하기로 약속합니다. 그런 다음, 이 증권을 외부 투자자에게 판매하여 수익을 얻습니다.
        \item \textbf{경제적 기능:}
        \begin{enumerate}
            \item \textbf{정보 수집 및 처리:} 잠재적 차입자에 대한 정보를 수집하고 처리하는 사업을 합니다. 자본 시장의 '문지기(gatekeepers)' 역할을 합니다.
            \item \textbf{저축자와 차입자 연결:} 인수 과정을 통해 저축자와 차입자를 연결합니다.
        \end{enumerate}
    \end{itemize}

    \item \textbf{투자 펀드 (Investment Funds):}
    \begin{itemize}
        \item \textbf{한 줄 핵심 요약:} 저축을 모아 대규모의 분산된 자산 포트폴리오에 투자하여 투자자들이 부를 저장하고 성장시킬 수 있도록 돕습니다.
        \item \textbf{종류:}
        \begin{itemize}
            \item \textbf{뮤추얼 펀드 (Mutual Funds) 및 상장지수펀드 (Exchange Traded Funds, ETFs):} 소액 투자자도 최소 투자 요건과 낮은 거래 수수료로 접근할 수 있습니다. 다양한 투자 프로필을 가집니다 (예: 머니마켓 펀드는 매우 안전하고 단기 채권에만 투자, 주식 뮤추얼 펀드는 더 위험). Vanguard와 같은 대형 뮤추얼 펀드 회사는 S\&P 500을 추종하는 펀드를 3,000달러의 최소 투자금으로 제공하여, 개인이 직접 이러한 투자를 복제하기 어려운 '풀링(pooling)'의 이점을 제공합니다.
            \item \textbf{헤지 펀드 (Hedge Funds):} 부유하고 정교한 투자자만을 대상으로 합니다. 100만 달러 이상의 순자산을 가진 99명 이하의 투자자 또는 500만 달러 이상의 순자산을 가진 499명 이하의 투자자를 대상으로 합니다. 일반 소액 투자자를 피함으로써 특정 규제를 우회할 수 있으며, 증권 및 파생상품 시장에서 고위험 고수익 투자 전략을 추구할 수 있습니다.
        \end{itemize}
    \end{itemize}
\end{itemize}

\textbf{4. 할부금융 회사 (Finance Companies)}
\begin{itemize}
    \item \textbf{한 줄 핵심 요약:} 은행과 유사하지만 특정 대출 시장에 특화된 금융 기관입니다.
    \item \textbf{주요 활동:}
    \begin{itemize}
        \item \textbf{소비자 신용 (Consumer Credit):} 가계 구매를 위한 신용을 제공합니다 (예: 주택 모기지, 가전제품 할부 대출, 자동차 대출). Ford Motor Credit Company는 자동차 판매 금융을 담당합니다.
        \item \textbf{사업 대출 (Business Lending):} 항공기 구매 금융부터 농업 장비 리스 제공까지 다양합니다.
    \end{itemize}
    \item \textbf{은행과의 유사점:} 대출 및 리스 제공이 주 활동이며, 차입자를 심사하고 모니터링하여 정보 비용을 최소화하는 등 은행과 유사하게 작동합니다.
    \item \textbf{은행과의 차이점 (자금 조달):} 예금으로 자금을 조달하지 않습니다. 대신 채권 발행이나 은행 대출을 통해 자금을 조달합니다.
\end{itemize}

\textbf{5. 정부 후원 기업 (Government Sponsored Enterprises, GSEs)}
\begin{itemize}
    \item \textbf{한 줄 핵심 요약:} 미국 모기지 시장의 핵심 주체로, 정부가 설립했으나 민간 형태로 운영되며 주택 소유를 지원합니다.
    \item \textbf{주요 기관:}
    \begin{itemize}
        \item \textbf{Fannie Mae (Federal National Mortgage Association):} 1930년대 설립.
        \item \textbf{Freddie Mac (Federal Home Loan Mortgage Corporation):} 1968년 설립.
        \item \textbf{Ginnie Mae (Government National Mortgage Corporation):} 1968년 설립. (정부가 전액 소유)
    \end{itemize}
    \item \textbf{역할:} 저소득층 및 소외 계층을 포함한 미국 내 주택 소유를 지원합니다. 이는 모기지 이자율을 보조함으로써 이루어집니다.
    \item \textbf{보조 방식:}
    \begin{itemize}
        \item \textbf{직접 대출:} 시장 금리보다 낮은 금리로 직접 대출을 제공합니다.
        \item \textbf{상환 보증:} 다른 금융 기관이 발행한 특정 부동산 금융 상품에 대한 상환을 보증합니다. 이러한 보증은 보험과 유사하게 작용하여 보조금 대출의 위험을 줄입니다.
    \end{itemize}
    \item \textbf{자금 조달 및 운용:} GSE는 채권 발행 또는 다른 금융 기관에 대한 대출 보증 판매를 통해 자금을 조달합니다. 이 자금은 주로 모기지 대출을 제공하고, 대출 보증과 관련된 보험금을 지급하는 데 사용됩니다.
\end{itemize}

\textbf{미국 금융 기관의 상대적 중요성 (1970년 vs. 2018년)}
미국 연방준비제도(Federal Reserve)의 자금 흐름(flow of funds) 데이터를 통해 은행과 주요 비은행 금융 기관의 자산 보유 추정치를 비교할 수 있습니다.

\begin{tcolorbox}[summarybox]
  \textbf{주요 시사점:}
  \begin{enumerate}
      \item \textbf{은행의 비중 감소:} 1970년에는 은행이 전체 자산의 절반 이상을 차지했지만, 2018년에는 25\% 미만으로 감소했습니다. 보험 회사도 비중이 감소했습니다.
      \item \textbf{뮤추얼 펀드의 부상:} 뮤추얼 펀드의 자산 비중은 5\% 미만에서 30\% 이상으로 크게 증가했습니다. 이는 Vanguard와 같은 투자 펀드가 제공하는 저비용 인덱스 투자 상품에 대한 접근성이 증가했음을 반영합니다.
  \end{enumerate}
\end{tcolorbox}

\begin{tcolorbox}[summarybox]
  \textbf{결론:} 금융 기관은 저축자로부터 차입자에게 자금을 전달하며, 이 과정에는 많은 은행 및 비은행 금융 기관이 참여합니다. 각 기관은 고유한 경제적 기능과 대차대조표를 가집니다. 전 세계적으로 은행은 여전히 중요한 역할을 하지만, 미국에서는 최근 몇 년간 비은행 금융 기관, 특히 뮤추얼 펀드의 중요성이 커졌습니다. 이러한 자금 재배분이 경제에 미치는 잠재적 영향(긍정적 또는 부정적)은 추후 논의될 주제입니다.
\end{tcolorbox}

\newpage
\subsection{3. 금융과 실물 경제 (Finance and the Real Economy)}
\addcontentsline{toc}{subsection}{3. 금융과 실물 경제}
이 섹션에서는 금융 시스템의 발전이 실물 경제 활동에 어떤 영향을 미치는지 살펴봅니다.

\subsubsection*{3.1. 금융 발전과 경제 활동 (Financial Development and Economic Activity)}
\addcontentsline{toc}{subsubsection}{3.1. 금융 발전과 경제 활동}

\begin{tcolorbox}[summarybox]
  \textbf{핵심 요약:} 금융 발전은 돈과 금융 시스템의 효율성과 품질 향상을 의미하며, 이는 경제 성장과 위험 분산에 큰 이점을 제공하지만, 동시에 금융 불안정성과 경제 위기를 초래할 수도 있습니다.
\end{tcolorbox}

\textbf{금융 발전의 의미}
금융 발전(Financial Development)은 돈과 금융 시스템이 더 빠르고 더 효율적으로 작동하게 되는 것을 의미합니다. 이는 금융 시스템이 자원을 가장 생산적인 투자 기회로 더 잘 배분할 수 있게 됨을 뜻합니다.

\textbf{금융 발전의 잠재적 이점:}
\begin{itemize}
    \item \textbf{경제 성장 (Economic Growth):} 효율적인 금융 시스템은 자본이 필요한 기업과 개인에게 자금을 효과적으로 공급하여 생산성을 높이고 경제 전반의 성장을 촉진합니다.
    \item \textbf{위험 분산 (Risk-Sharing):} 보험, 뮤추얼 펀드와 같은 금융 상품을 통해 개인과 기업이 위험을 분산하고 관리할 수 있게 합니다. 이는 불확실성을 줄이고 투자를 장려합니다.
\end{itemize}

\textbf{금융 발전의 잠재적 비용 및 위험:}
\begin{itemize}
    \item \textbf{금융 불안정성 (Financial Instability):} 금융 시스템이 너무 빠르게 발전하거나 복잡해지면, 규제가 미치지 못하는 영역이 생기거나 새로운 유형의 위험이 발생할 수 있습니다. 이는 시스템 전체의 불안정성을 초래할 수 있습니다.
    \item \textbf{경제 위기 (Economic Crashes):} 금융 불안정성은 금융 위기로 이어질 수 있으며, 이는 실물 경제에 심각한 타격을 주어 경기 침체나 경제 위기를 유발할 수 있습니다. (예: 2008년 글로벌 금융 위기)
\end{itemize}

\begin{tcolorbox}[summarybox]
  \textbf{모듈 1의 최종 목표:} 이 모듈을 통해 돈과 금융이 무엇인지, 그리고 그것이 가져올 수 있는 잠재적인 이점과 비용에 대한 큰 그림을 이해하는 것입니다. 금융 발전은 양날의 검과 같아서, 그 혜택을 누리면서도 위험을 관리하는 것이 중요합니다.
\end{tcolorbox}

\newpage
\section{빠르게 훑어보기 (1페이지 요약)}
\addcontentsline{toc}{section}{빠르게 훑어보기 (1페이지 요약)}

\begin{tcolorbox}[summarybox={모듈 1: 돈과 금융 핵심 요약}]
  \textbf{1. 돈의 본질과 기능:}
  \begin{itemize}
    \item 돈은 \textbf{교환의 매개}, \textbf{가치 척도}, \textbf{가치 저장 수단}의 세 가지 핵심 기능을 수행합니다.
    \item \textbf{상품 화폐} (본질적 가치) $\rightarrow$ \textbf{종이 화폐} (금속 보증) $\rightarrow$ \textbf{법정 화폐} (정부 신뢰)로 진화했습니다.
    \item \textbf{돈 $\neq$ 부}: 돈은 부의 한 형태일 뿐, 모든 자산이 돈은 아닙니다.
  \end{itemize}

  \textbf{2. 지불 시스템:}
  \begin{itemize}
    \item \textbf{현금}과 \textbf{비현금 결제} (수표, ACH, 카드, 디지털 지갑)로 나뉩니다.
    \item \textbf{전자 결제}가 증가 추세이며, 지불 시스템은 경제의 '금융 배관' 역할을 합니다.
    \item \textbf{비트코인}은 현재 돈의 핵심 기능(지불 수단, 가치 저장)을 충족하지 못합니다.
  \end{itemize}

  \textbf{3. 통화량과 인플레이션:}
  \begin{itemize}
    \item \textbf{M1} (현금 + 요구불 예금)과 \textbf{M2} (M1 + 저축성 예금)로 통화량을 측정합니다.
    \item \textbf{인플레이션}은 물가 상승률이며, \textbf{CPI}로 측정합니다.
    \item 과거에는 \textbf{통화량 증가}와 \textbf{인플레이션} 간의 강한 상관관계가 있었으나 (1970년대 대인플레이션), 21세기에는 그 관계가 약화되었습니다.
  \end{itemize}

  \textbf{4. 금융 시스템의 역할:}
  \begin{itemize}
    \item \textbf{저축자} (자금 여유)와 \textbf{차입자} (자금 수요)를 연결하여 자원을 효율적으로 배분합니다.
    \item \textbf{금융 상품} (예: 주식, 채권, 예금, 보험)을 통해 자금 교환이 이루어집니다.
  \end{itemize}

  \textbf{5. 금융 시장 vs. 금융 기관:}
  \begin{itemize}
    \item \textbf{금융 시장:} 브로커/딜러를 통해 직접 거래 (예: 주식 시장). 낮은 거래 비용, 효율적 정보, 투자자 보호가 중요합니다.
    \item \textbf{금융 기관:} 은행과 같은 중개자를 통해 간접 거래. \textbf{만기 불일치}, \textbf{가격 위험}, \textbf{정보 비대칭} 문제를 해결하여 가치를 창출합니다.
  \end{itemize}

  \textbf{6. 은행 vs. 비은행 금융 기관:}
  \begin{itemize}
    \item \textbf{은행:} 예금을 받고 대출을 해주는 전통적 기관 (상업 은행, 저축 기관 등).
    \item \textbf{비은행 금융 기관:} 예금을 받지 않으면서 다양한 금융 서비스 제공 (보험사, 연기금, 증권사, 할부금융사, 정부 후원 기업 등).
    \item 미국에서는 \textbf{은행의 비중이 감소}하고 \textbf{뮤추얼 펀드}와 같은 비은행 기관의 중요성이 증가하는 추세입니다.
  \end{itemize}

  \textbf{7. 금융과 실물 경제:}
  \begin{itemize}
    \item \textbf{금융 발전}은 경제 성장과 위험 분산에 기여하지만, 금융 불안정성과 경제 위기를 초래할 수도 있습니다.
  \end{itemize}
\end{tcolorbox}

\newpage

\newpage

\begin{tcolorbox}[title=\textbf{개요: 금융 발전과 경제 활동}]
금융과 화폐 시스템은 경제의 원활한 기능에 필수적입니다. 금융 발전은 사회 전반의 경제적 성과를 향상시키지만, 금융 시스템의 교란은 막대한 비용을 초래할 수 있습니다. 본 문서에서는 금융 발전이 경제 성장을 촉진하는 증거를 제시하고, 금융 위기의 발생 원인과 사회적 비용을 논의합니다. 궁극적으로 금융 효율성과 안정성 사이의 균형을 찾는 것이 중요하며, 이를 위한 정책적 질문들을 탐구합니다.
\end{tcolorbox}

\section{Lesson 1-3.1: 금융 발전과 경제 활동}

\begin{tcolorbox}[title=\textbf{용어 정리}]
\begin{adjustbox}{width=\textwidth,center}
\begin{tabular}{lllc}
\toprule
\textbf{용어} & \textbf{쉬운 설명} & \textbf{원어} & \textbf{비고} \\
\midrule
금융 발전 & 금융 시장과 기관이 성장하여 자금 흐름이 효율화되는 과정 & Financial Development & 경제 성장의 핵심 동력 \\
금융 시스템 & 저축자의 자금을 차입자에게 연결하는 시장과 기관의 총체 & Financial System & 은행, 비은행 금융기관, 금융 시장 포함 \\
금융 위기 & 금융 시스템 내에서 투자자 신뢰 상실 등으로 발생하는 심각한 교란 & Financial Crisis & 경제에 막대한 피해를 줄 수 있음 \\
증권화 & 위험한 대출을 묶어 새로운 유가증권을 만드는 과정 & Securitization & 2008년 금융 위기의 주요 원인 중 하나 \\
자산유동화증권 & 대출 채권 등 자산을 담보로 발행되는 증권 & Asset-Backed Securities (ABS) & 위험을 분산하거나 변환하는 데 사용 \\
서브프라임 모기지 & 신용 등급이 낮은 사람들에게 제공되는 주택 담보 대출 & Subprime Mortgages & 부실 위험이 높은 대출 \\
머니 마켓 펀드 & 단기 채권에 투자하는 안전 지향 뮤추얼 펀드 & Money Market Funds & 은행 예금과 유사하지만 정부 보증 없음 \\
초인플레이션 & 물가 상승률이 극도로 높아져 화폐 가치가 급락하는 현상 & Hyperinflation & 정부의 과도한 화폐 발행이 원인 \\
디지털 화폐 & 블록체인 기술 기반의 전자 화폐 (예: 비트코인) & Digital Currencies & 정부 통제 밖의 화폐 \\
스테이블코인 & 가치가 법정 화폐에 고정되도록 설계된 디지털 화폐 & Stablecoins & 테더(Tether)가 대표적 \\
상업어음 & 기업이 단기 자금을 조달하기 위해 발행하는 무담보 어음 & Commercial Paper & 단기 금융 시장의 중요한 도구 \\
정부 구제금융 & 금융 위기 시 정부가 금융 기관에 제공하는 재정 지원 & Government Bailout & 시스템 붕괴 방지를 위한 최후의 수단 \\
\bottomrule
\end{tabular}
\end{adjustbox}
\end{tcolorbox}

\newpage

\section{핵심 개념 및 원리}

\subsection{금융 발전과 경제 활동의 관계}

\begin{tcolorbox}[title=\textbf{핵심 요약}]
금융 발전은 경제 성장을 촉진하고 사회 전반의 경제적 성과를 향상시키는 핵심 동력입니다. 이는 자금의 효율적인 배분, 시장 가격의 정확성, 그리고 위험 공유의 개선을 통해 이루어집니다.
\end{tcolorbox}

\subsubsection*{금융 시스템의 역할}
경제의 원활한 기능은 화폐와 지급 시스템에 크게 의존합니다. 금융 시스템은 금융 시장과 다양한 은행 및 비은행 금융 기관을 포함하며, 저축자(savers)로부터 차입자(borrowers)에게 자금을 효율적으로 전달하는 중요한 역할을 합니다. 이러한 자금의 흐름은 투자를 촉진하고, 이는 다시 경제 성장의 기반이 됩니다.

\begin{itemize}
    \item \textbf{자금 중개 (Financial Intermediation):} 금융 시스템은 저축된 자금을 필요로 하는 기업이나 개인에게 연결하여 투자를 가능하게 합니다.
    \item \textbf{자원 배분 효율성 (Efficient Resource Allocation):} 금융 발전은 시장 가격을 더욱 정확하게 만들고, 자원이 가장 생산적인 곳으로 배분되도록 돕습니다.
    \item \textbf{위험 공유 (Risk Sharing):} 다양한 금융 상품과 시장을 통해 투자 위험을 여러 주체에게 분산시켜 개별 경제 주체의 부담을 줄여줍니다.
\end{itemize}
이러한 요소들이 결합될 때, 우리는 더 빠른 경제 성장과 사회 전반의 더 나은 경제적 성과를 기대할 수 있습니다.

\subsubsection*{금융 발전의 측정 및 증거}
금융 발전이 경제 성장에 미치는 긍정적인 영향을 확인하기 위해, 우리는 세계은행(World Bank)의 최신 데이터를 활용하여 이 아이디어를 검증할 수 있습니다.

\begin{itemize}
    \item \textbf{측정 지표:} 금융 발전을 측정하기 위해 'GDP 대비 민간 부문 신용 총액(total amount of credit extended to the private sector scaled by GDP)'을 사용합니다. 이는 국가 간 비교를 위해 GDP로 신용 규모를 표준화한 것입니다.
    \item \textbf{경제 발전 지표:} 경제 발전의 척도로는 '1인당 GDP(GDP per capita)'를 사용합니다.
    \item \textbf{분석 결과:} 전 세계 국가들을 대상으로 한 산점도(scatter plot) 분석 결과는 우리의 직관을 뒷받침합니다. 금융 시스템 규모가 큰 국가일수록 시민들의 소득(1인당 GDP)이 더 높은 경향을 보입니다.
\end{itemize}
이는 금융 발전이 긍정적인 경제적 결과를 가져온다는 강력한 증거입니다. 따라서 정책 입안자라면 더 많은 금융 발전, 더 많은 은행, 더 많은 금융 시장, 그리고 이러한 기관에 대한 규제 완화를 우선시해야 한다고 생각할 수 있습니다.

\begin{tcolorbox}[colback=lightgreen, colframe=green!75!black, title=\textbf{정책적 시사점}] % mygreen -> lightgreen
금융 발전이 경제 성장에 긍정적인 영향을 미친다는 증거는 정책 입안자들이 금융 부문의 성장을 장려해야 함을 시사합니다. 이는 금융 시장의 효율성을 높이고, 자금 조달을 용이하게 하며, 궁극적으로 국가 경제의 번영에 기여할 수 있습니다.
\end{tcolorbox}

\subsection{금융 시스템의 교란과 금융 위기}
금융 발전이 항상 순조로운 것만은 아닙니다. 때로는 금융 시스템 내에서 심각한 교란이 발생할 수 있으며, 이는 사회에 막대한 비용을 초래하는 금융 위기로 이어집니다.

\subsubsection*{금융 위기의 발생 원인}
금융 위기는 주로 두 가지 주요 원인으로 인해 발생합니다.
\begin{enumerate}
    \item \textbf{투자자 신뢰 상실:} 투자자들이 금융 기관에 대한 신뢰를 잃을 때 발생합니다. 이는 은행 예금 인출 사태(bank run)와 같은 형태로 나타날 수 있습니다.
    \item \textbf{빠르고 통제되지 않는 금융 혁신:} 금융 혁신이 너무 빠르게 진행되고 시장이나 규제 당국에 의해 제대로 검증되지 않을 때 발생합니다. 새로운 금융 상품이나 기술의 위험이 충분히 평가되지 않아 시스템 전체의 불안정성을 초래할 수 있습니다.
\end{enumerate}

\subsubsection*{금융 위기의 전개 과정}
금융 위기는 일반적으로 다음과 같은 단계로 전개되며, 이는 대공황(Great Depression)이나 대침체(Great Recession)와 같은 역사적 사건에서 반복적으로 관찰된 패턴입니다.
\begin{enumerate}
    \item \textbf{과잉과 공황 (Excesses and Panic):} 금융 시스템 내의 과도한 위험 추구나 거품이 공황을 유발합니다.
    \item \textbf{자금 인출 (Withdrawals):} 투자자들이 금융 기관에서 저축을 대규모로 인출하기 시작합니다.
    \item \textbf{자산 청산 (Liquidation):} 자금 인출에 대응하기 위해 금융 기관들은 투자 자산을 강제로 청산해야 합니다. 이는 자산 가격 하락을 더욱 부추깁니다.
    \item \textbf{시장 붕괴 (Market Crash):} 주식 시장과 같은 금융 시장이 붕괴됩니다.
    \item \textbf{경제적 부담 (Economic Burden):} 가계는 저축을 잃고, 신용에 접근할 수 없게 됩니다. 기업들은 자금 조달에 어려움을 겪어 파산합니다.
    \item \textbf{장기 불황 (Deep and Long-lasting Recession):} 이 모든 것이 깊고 장기적인 경기 침체로 이어집니다.
\end{enumerate}
이러한 위기 이야기의 초기 버전에서는 은행이 중심적인 역할을 했습니다.

\subsubsection*{역사 속의 금융 공황}
미국에서는 1793년부터 1797년, 1811년, 1813년, 1816년, 그리고 1873년, 1894년, 1890년, 1893년 등 수많은 공황(panics)이 발생했습니다. 이러한 기록을 가진 시스템의 진정한 이름은 '공황 시스템(panic system)'이라고 불릴 정도였습니다.

\begin{itemize}
    \item \textbf{1907년 공황 (Panic of 1907):} 20세기 최초의 세계적인 금융 위기였습니다.
    \item \textbf{1930년대 대공황 (Great Depression):} 1930년대 대공황 기간 동안 은행 공황은 절정에 달했습니다. 1930년에서 1933년 사이에 매년 약 2,000개의 은행이 파산했습니다. 이는 현대 역사상 최악의 글로벌 금융 위기였을 것입니다.
    \item \textbf{FDIC 설립:} 대공황의 경험은 현재의 정부 예금 보험 시스템인 연방예금보험공사(Federal Deposit Insurance Corporation, FDIC)의 설립으로 이어졌습니다. 이는 예금자들을 보호하고 은행 시스템의 안정성을 높이기 위한 조치였습니다.
\end{itemize}
100년이 지난 지금도 금융 위기에는 여전히 공황에 빠진 투자자들이 관련되어 있습니다.

\newpage

\subsection{현대 금융 위기의 주요 원인}
21세기 금융 위기는 과거와는 다른 새로운 형태의 금융 혁신과 시장 구조에서 비롯되기도 합니다.

\subsubsection*{증권화 (Securitization)}
증권화는 20세기 위대한 금융 혁신 중 하나입니다.
\begin{itemize}
    \item \textbf{개념:} 금융 공학자들이 위험한 대출 풀(pools of risky loans)을 모아 새로운 자산유동화증권(Asset-Backed Securities, ABS)을 만드는 것을 가능하게 합니다.
    \item \textbf{장점 (겉보기):} 일부 ABS는 트리플 A(Triple A) 등급을 받아, 위험을 안전으로 변환하는 것처럼 보였습니다.
    \item \textbf{문제점:} 그러나 증권화 과정 자체는 복잡하고 불투명하여 많은 투자자들이 이해하기 어려웠습니다.
    \item \textbf{2007년 위기 사례:} 2007년부터 가장 위험한 모기지 담보 증권(mortgage-backed securities), 특히 서브프라임 모기지(subprime mortgages)가 부실화되기 시작했습니다. 투자자들은 이것이 불가능하다고 생각했기 때문에 공황에 빠졌습니다. 이 공황은 모기지와 전혀 관련 없는 학자금 대출 및 신용카드 증권화를 포함한 모든 증권화 시장으로 확산되어 신용 시장의 거의 완전한 붕괴를 초래했습니다.
\end{itemize}

\subsubsection*{머니 마켓 펀드 (Money Market Funds)}
머니 마켓 펀드는 21세기 금융 공황의 중심에 있었습니다.
\begin{itemize}
    \item \textbf{개념:} 이 펀드들은 단기 채권 시장에 매우 안전한 투자를 하는 뮤추얼 펀드의 일종입니다. 대규모 기관 투자자들은 이 펀드에 현금을 예치하여 약간의 수익을 얻습니다.
    \item \textbf{특징:} 이 펀드들은 부채 구조를 은행 예금과 유사하게 만들어 위험이 없어 보이고 요구 시 상환이 가능하도록 합니다.
    \item \textbf{위험:} 그러나 머니 마켓 펀드의 저축은 정부에 의해 보증되지 않으며 손실을 입을 수 있습니다.
    \item \textbf{2008년 위기 사례:} 2008년, 리저브 프라이머리 펀드(Reserve Primary Fund)라는 유명한 펀드에서 손실이 발생했습니다. 이는 전혀 예상치 못한 일이었고, 투자자들은 공황에 빠져 전체 머니 마켓 펀드 산업에서 대규모 자금 인출이 발생하여 붕괴 직전까지 몰렸습니다.
\end{itemize}

\subsubsection*{화폐 및 지급 시스템 문제}
화폐와 지급 시스템 자체에서도 문제가 발생할 수 있습니다.
\begin{itemize}
    \item \textbf{초인플레이션 (Hyperinflation):} 1920년대 바이마르 독일, 1946년 헝가리, 2008년 짐바브웨, 2018년 베네수엘라의 초인플레이션은 극단적인 예시입니다.
    \item \textbf{원인:} 각 사례에서 정부는 너무 많은 지폐를 발행했고, 결국 지폐는 가치를 상실했습니다.
    \item \textbf{결과:} 이는 경제와 사회 전체에 참으로 끔찍한 결과를 초래했습니다.
\end{itemize}

\newpage

\subsection{디지털 화폐와 금융 시스템}
디지털 화폐의 등장은 화폐 시스템에 새로운 도전과 기회를 동시에 제공하고 있습니다.

\subsubsection*{비트코인 (Bitcoin)}
\begin{itemize}
    \item \textbf{특징:} 디지털 화폐는 정부의 지폐 공급 통제를 벗어납니다. 2008년 비트코인 등장 이후, 비트코인은 신뢰할 수 있는 지급 수단이라기보다는 투기적 거품처럼 보였습니다.
    \item \textbf{위험 평가:} 한 가지 관점은 가격 위험이 소매 투자자들 사이에 분산되어 있는 한, 디지털 화폐의 가치가 오르든 내리든 크게 중요하지 않다는 것입니다. 이는 투자자들의 문제일 뿐, 금융 시스템이나 광범위한 경제에 대한 위협은 아니라는 주장입니다.
\end{itemize}

\subsubsection*{스테이블코인 (Stablecoins)}
스테이블코인은 비트코인과는 다릅니다.
\begin{itemize}
    \item \textbf{개념:} 스테이블코인은 가치를 일반 화폐(예: 달러)에 고정시키는 특별한 유형의 디지털 화폐입니다.
    \item \textbf{테더 (Tether) 사례:} 테더는 가장 인기 있는 스테이블코인입니다. 2021년 기준으로 테더는 600억 달러 이상의 디지털 토큰을 발행했으며, 이 토큰들은 개당 1달러로 교환될 수 있습니다. 이 스테이블코인들은 물건을 구매하는 데 사용될 수 있으며, 달러로 교환될 수 있기 때문에 가치를 가집니다.
    \item \textbf{숨겨진 위험:} 테더는 이러한 디지털 토큰을 판매하여 모금한 자금을 투자합니다. 알려진 바에 따르면, 테더는 모은 달러를 상업어음(commercial paper)과 같은 유가증권에 투자합니다. 2021년 기준으로 테더는 약 300억 달러 상당의 상업어음 투자를 보유하고 있었으며, 이는 사실상 세계에서 가장 큰 머니 마켓 펀드 중 하나였습니다.
    \item \textbf{위기 시나리오:} 만약 투자자들이 테더에 대한 신뢰를 잃고 대규모로 디지털 토큰을 상환(redeem)한다면 어떻게 될까요? 일부 투자자들이 공황에 빠져 토큰당 1달러를 받지 못할까 걱정한다면, 그들은 90센트 또는 그 이하의 가격이라도 기꺼이 받아들일 수 있습니다. 이는 테더에 대한 뱅크런(run)으로 이어질 수 있으며, 디지털 화폐의 가치는 즉시 폭락할 수 있습니다. 이러한 공황은 전체 스테이블코인 시장을 교란시킬 수 있습니다.
    \item \textbf{파급 효과:} 더 나쁜 것은, 이러한 교란이 상업어음 시장으로 파급될 수 있다는 점입니다. 이는 정부와 기업이 운전자본을 조달하는 방식에 영향을 미칠 수 있습니다. 일부 디지털 화폐는 다른 화폐보다 더 문제가 될 수 있음을 시사합니다.
\end{itemize}

\subsection{정부 구제금융 (Government Bailouts)}
금융 시스템이 실패할 때마다 정부가 개입하여 지원을 제공하는 것을 우리는 반복적으로 목격했습니다. 이러한 조치를 일반적으로 '정부 구제금융(government bailout)'이라고 합니다.

\begin{itemize}
    \item \textbf{목적:} 금융 시스템의 붕괴를 막고, 그로 인한 경제적 피해를 최소화하는 것이 주된 목적입니다.
    \item \textbf{2008년 미국 사례:} 미국에서는 2008년 금융 위기 동안 및 이후에 다양한 정부 기관이 6천억 달러 이상의 납세자 자금을 지출했습니다. 이는 막대한 금액이지만, 미국은 매우 큰 경제 규모를 가지고 있어 개입이 성공한다면 그 정도의 돈을 지출할 여력이 있습니다. 다행히 2008년에는 대체로 성공적으로 작동했습니다.
    \item \textbf{다른 국가의 한계:} 그러나 다른 정부들은 그 정도 규모의 구제금융을 감당할 수 없습니다.
    \item \textbf{비용의 상대적 규모:} 최근 역사상 가장 큰 은행 구제금융 사례를 GDP 대비 규모로 보면, 1997년 인도네시아는 국가 소득의 절반 이상을 은행 구제금융에 지출했습니다. 이 정도 규모의 구제금융은 정부 재정의 갑작스러운 붕괴로 이어질 수 있습니다.
\end{itemize}

\newpage

\subsection{금융 정책의 트레이드오프와 규제 질문}
금융 발전과 안정성 사이에는 본질적인 트레이드오프(tradeoff)가 존재합니다.

\begin{tcolorbox}[title=\textbf{핵심 트레이드오프}]
우리는 효율적이고 경쟁적인 금융 부문을 육성하고, 화폐와 금융의 혁신이 경제를 발전시키도록 장려해야 합니다. 그러나 동시에 금융 부문의 안정성과 위기 발생 위험 사이에서 균형을 찾아야 합니다.
\end{tcolorbox}

위기를 방지하기 위해서는 한계를 설정해야 합니다. 이는 중요한 공공 정책 질문으로 이어집니다.

\begin{enumerate}
    \item \textbf{규제 대상 기관:} 어떤 기관을 규제해야 하는가? 사모펀드(private equity)와 핀테크(FinTech) 기업은 규제 당국이 그대로 두어야 하는가? 이는 여전히 열린 질문입니다.
    \item \textbf{규제 기관의 허용 범위:} 규제 대상 금융 기관은 무엇을 할 수 있도록 허용되어야 하는가?
    \begin{itemize}
        \item 테더는 은행과 동일한 공시 의무를 가져야 하는가?
        \item 은행은 헤지펀드 자회사(hedge funds subsidiaries)를 가질 수 있도록 허용되어야 하는가?
        \item 이러한 헤지펀드 자회사가 비트코인에 투자하는 것을 허용해야 하는가?
    \end{itemize}
    \item \textbf{금융 위기 시 정부 개입:} 금융 위기 동안 어떤 종류의 정부 개입이 적절한가? 우리의 직관은 때때로 실패하는 기관을 구제하는 것이 적절하다는 것입니다. 그러나 그 한계는 어떻게 설정되어야 하는가?
\end{enumerate}
이러한 질문들은 금융 시스템의 미래와 경제 전반의 안정성을 결정하는 데 매우 중요합니다.

\newpage

\begin{tcolorbox}[title=\textbf{빠르게 훑어보기: Module 1 핵심 요약}]
\begin{itemize}
    \item \textbf{금융 발전의 중요성:} 금융 시스템의 발전은 자금 중개, 자원 배분 효율성, 위험 공유를 통해 경제 성장을 촉진합니다. (예: GDP 대비 민간 신용 규모가 클수록 1인당 GDP 높음)
    \item \textbf{금융 위기의 그림자:} 금융 시스템은 투자자 신뢰 상실이나 통제되지 않는 혁신으로 교란될 수 있으며, 이는 공황, 자금 인출, 시장 붕괴, 장기 불황으로 이어지는 금융 위기를 초래합니다. (예: 1930년대 대공황, 2008년 대침체)
    \item \textbf{현대 위기의 원인:} 증권화(서브프라임 모기지), 머니 마켓 펀드(2008년 리저브 프라이머리 펀드), 초인플레이션(정부의 과도한 화폐 발행) 등이 현대 금융 위기의 주요 원인입니다.
    \item \textbf{디지털 화폐의 양면성:} 비트코인은 투기적 성격이 강하며, 스테이블코인(테더)은 안정성을 추구하지만, 대규모 인출 시 상업어음 시장 등 광범위한 금융 시스템에 위험을 전이시킬 수 있습니다.
    \item \textbf{정부의 역할과 딜레마:} 금융 위기 시 정부 구제금융은 시스템 붕괴를 막지만, 막대한 비용과 정부 재정 부담을 초래할 수 있습니다. (예: 1997년 인도네시아 GDP의 절반 이상 지출)
    \item \textbf{정책적 트레이드오프:} 금융 효율성과 혁신을 장려하는 동시에 금융 시스템의 안정성을 확보하는 것이 중요합니다. 규제 대상, 규제 범위, 위기 시 개입 한계에 대한 지속적인 논의가 필요합니다.
\end{itemize}
\end{tcolorbox}

\end{document}
